\input{preamble}

% OK, start here.
%
\begin{document}

\title{Hyperrecouvrements}


\maketitle

\phantomsection
\label{section-phantom}

\tableofcontents

\section{Introduction}
\label{section-introduction}

\noindent
Soit $\mathcal{C}$ un site, voir Sites, définition \ref{sites-definition-site}.
Soit $X$ un objet de $\mathcal{C}$.
Étant donné un faisceau abélien $\mathcal{F}$
sur $\mathcal{C}$, nous voudrions calculer
ses groupes de cohomologie
$$
H^i(X, \mathcal{F}).
$$
Selon nos définitions générales (Cohomologie sur les sites, section
\ref{sites-cohomology-section-cohomology-sheaves})
ce groupe de cohomologie se calcule en
choisissant une résolution injective
$
0 \to \mathcal{F} \to \mathcal{I}^0 \to \mathcal{I}^1 \to \ldots
$
et en posant
$$
H^i(X, \mathcal{F})
=
H^i(
\Gamma(X, \mathcal{I}^0) \to
\Gamma(X, \mathcal{I}^1) \to
\Gamma(X, \mathcal{I}^2)\to \ldots)
$$
Le but de ce chapitre est de montrer que nous pouvons aussi calculer ces
groupes de cohomologie sans choisir de résolution injective
(dans le cas où $\mathcal{C}$ admet des produits fibrés). Pour ce faire,
nous utiliserons des hyperrecouvrements.

\medskip\noindent
Un hyperrecouvrement dans un site est une généralisation d'un recouvrement, voir
\cite[Expos\'e V, Sec. 7]{SGA4}. Étant donné un hyperrecouvrement $K$ d'un objet
$X$, il existe une suite spectrale de {\v C}ech vers la cohomologie
qui exprime la cohomologie d'un faisceau abélien $\mathcal{F}$
sur $X$ en fonction de la cohomologie du faisceau sur les
composantes $K_n$ de $K$. Il s'avère qu'il existe toujours
assez d'hyperrecouvrements pour que, en prenant la limite inductive sur tous les hyperrecouvrements,
la suite spectrale dégénère et que la cohomologie de $\mathcal{F}$
sur $X$ soit calculée par la limite inductive des groupes de cohomologie de {\v C}ech.

\medskip\noindent
Un objet plus général que l'on peut considérer est une augmentation simpliciale pour laquelle
on a la descente cohomologique, voir \cite[Expos\'e Vbis]{SGA4}. Un excellent
manuscrit sur la descente cohomologique est le texte de Brian Conrad, voir
\url{https://math.stanford.edu/~conrad/papers/hypercover.pdf}.
Nous reviendrons sur ces questions dans le chapitre sur les espaces simpliciaux,
où nous montrerons, par exemple, que les hyperrecouvrements propres d'espaces
topologiques « localement compacts » sont de descente
cohomologique (Espaces simpliciaux, section
\ref{spaces-simplicial-section-proper-hypercovering}).
Notre méthode consistera à ramener cet énoncé à la suite spectrale de {\v C}ech vers
la cohomologie construite dans ce chapitre.






























\section{Objets semi-représentables}
\label{section-semi-representable}

\noindent
Pour commencer, nous posons la définition suivante.
Les lettres « SR » signifient semi-représentable.

\begin{definition}
\label{definition-SR}
Soit $\mathcal{C}$ une catégorie. Nous notons $\text{SR}(\mathcal{C})$
la catégorie des {\it objets semi-représentables} définie comme suit :
\begin{enumerate}
\item les objets sont les familles d'objets $\{U_i\}_{i \in I}$ ;
\item les morphismes $\{U_i\}_{i \in I} \to \{V_j\}_{j \in J}$ sont donnés par
une application $\alpha : I \to J$ et, pour chaque $i \in I$,
un morphisme $f_i : U_i \to V_{\alpha(i)}$ de $\mathcal{C}$.
\end{enumerate}
Soit $X \in \Ob(\mathcal{C})$ un objet de $\mathcal{C}$.
La catégorie des {\it objets semi-représentables au-dessus de $X$}
est la catégorie
$\text{SR}(\mathcal{C}, X) = \text{SR}(\mathcal{C}/X)$.
\end{definition}

\noindent
Cette définition est essentiellement équivalente à
\cite[Expos\'e V, Subsection 7.3.0]{SGA4}. Notons qu'il
s'agit d'une « grande » catégorie. Nous bornerons plus loin la taille des ensembles
d'indices $I$ dont nous avons besoin pour les hyperrecouvrements de $X$. Nous pourrons alors redéfinir
$\text{SR}(\mathcal{C}, X)$ de façon à en faire une catégorie. Explicitons
les objets et les morphismes de $\text{SR}(\mathcal{C}, X)$ :
\begin{enumerate}
\item les objets sont les familles de morphismes
$\{U_i \to X\}_{i \in I}$ ;
\item les morphismes $\{U_i \to X\}_{i \in I} \to
\{V_j \to X\}_{j \in J}$ sont donnés par
une application $\alpha : I \to J$ et, pour chaque $i \in I$,
un morphisme $f_i : U_i \to V_{\alpha(i)}$ au-dessus de $X$.
\end{enumerate}
Il existe un foncteur d'oubli
$\text{SR}(\mathcal{C}, X) \to \text{SR}(\mathcal{C})$.

\begin{definition}
\label{definition-SR-F}
Soit $\mathcal{C}$ une catégorie.
Nous notons $F$ le foncteur {\it qui associe un préfaisceau à un
objet semi-représentable}. En formule,
\begin{eqnarray*}
F : \text{SR}(\mathcal{C}) & \longrightarrow & \textit{PSh}(\mathcal{C}) \\
\{U_i\}_{i \in I} & \longmapsto & \amalg_{i\in I} h_{U_i}
\end{eqnarray*}
où $h_U$ désigne le préfaisceau représentable associé à
l'objet $U$.
\end{definition}

\noindent
Étant donné un morphisme $U \to X$, nous obtenons un morphisme $h_U \to h_X$ de préfaisceaux
représentables. Ainsi, nous considérons souvent la restriction de $F$ à $\text{SR}(\mathcal{C}, X)$
comme un foncteur à valeurs dans la catégorie des préfaisceaux d'ensembles au-dessus de $h_X$,
à savoir $\textit{PSh}(\mathcal{C})/h_X$. Voici un diagramme :
$$
\xymatrix{
\text{SR}(\mathcal{C}, X) \ar[r]_F \ar[d] &
\textit{PSh}(\mathcal{C})/h_X \ar[d] \\
\text{SR}(\mathcal{C}) \ar[r]^F &
\textit{PSh}(\mathcal{C})
}
$$
Examinons maintenant l'existence de limites dans la catégorie des objets
semi-représentables.

\begin{lemma}
\label{lemma-coprod-prod-SR}
Soit $\mathcal{C}$ une catégorie.
\begin{enumerate}
\item la catégorie $\text{SR}(\mathcal{C})$ admet des coproduits
et $F$ commute à ceux-ci ;
\item le foncteur $F : \text{SR}(\mathcal{C}) \to \textit{PSh}(\mathcal{C})$
commute aux limites ;
\item si $\mathcal{C}$ admet des produits fibrés, alors $\text{SR}(\mathcal{C})$
admet des produits fibrés ;
\item si $\mathcal{C}$ admet les produits de deux objets, alors
$\text{SR}(\mathcal{C})$ admet les produits de deux objets ;
\item si $\mathcal{C}$ admet des égalisateurs, il en va de même de $\text{SR}(\mathcal{C})$ ;
\item si $\mathcal{C}$ admet un objet final, il en va de même de $\text{SR}(\mathcal{C})$.
\end{enumerate}
Soit $X \in \Ob(\mathcal{C})$.
\begin{enumerate}
\item la catégorie $\text{SR}(\mathcal{C}, X)$ admet des coproduits
et $F$ commute à ceux-ci ;
\item si $\mathcal{C}$ admet des produits fibrés, alors $\text{SR}(\mathcal{C}, X)$
admet des limites finies et
$F : \text{SR}(\mathcal{C}, X) \to \textit{PSh}(\mathcal{C})/h_X$
commute à celles-ci.
\end{enumerate}
\end{lemma}

\begin{proof}
Démonstration des résultats sur $\text{SR}(\mathcal{C})$.
Démonstration de (1). Le coproduit de $\{U_i\}_{i \in I}$ et $\{V_j\}_{j \in J}$ est
$\{U_i\}_{i \in I} \amalg \{V_j\}_{j \in J}$, autrement dit la famille
d'objets dont l'ensemble d'indices est $I \amalg J$ et qui, à un élément
$k \in I \amalg J$, associe $U_i$ si $k = i \in I$ et $V_j$ si
$k = j \in J$. Il en va de même pour les coproduits
de familles d'objets. Il est clair que $F$ commute à ceux-ci.

\medskip\noindent
Démonstration de (2). Pour $U$ dans $\Ob(\mathcal{C})$, considérons l'objet $\{U\}$ de
$\text{SR}(\mathcal{C})$. Il est clair que
$\Mor_{\text{SR}(\mathcal{C})}(\{U\}, K)) = F(K)(U)$
pour $K \in \Ob(\text{SR}(\mathcal{C}))$. Puisque les limites de préfaisceaux
se calculent au niveau des sections
(Sites, section \ref{sites-section-limits-colimits-PSh}),
nous concluons que $F$ commute aux limites.

\medskip\noindent
Démonstration de (3). Supposons donnés un morphisme
$(\alpha, f_i) : \{U_i\}_{i \in I} \to \{V_j\}_{j \in J}$
et un morphisme
$(\beta, g_k) : \{W_k\}_{k \in K} \to \{V_j\}_{j \in J}$.
Le produit fibré de ces morphismes est donné par
$$
\{ U_i \times_{f_i, V_j, g_k} W_k\}_{(i, j, k) \in I \times J \times K
\text{ tels que } j = \alpha(i) = \beta(k)}
$$
Les produits fibrés existent si $\mathcal{C}$ admet des produits fibrés.

\medskip\noindent
Démonstration de (4). Le produit de $\{U_i\}_{i \in I}$ et $\{V_j\}_{j \in J}$ est
$\{U_i \times V_j\}_{i \in I, j \in J}$. Les produits existent si
$\mathcal{C}$ admet des produits.

\medskip\noindent
Démonstration de (5). L'égalisateur de deux applications
$(\alpha, f_i), (\alpha', f'_i) : \{U_i\}_{i \in I} \to \{V_j\}_{j \in J}$
est
$$
\{
\text{Eq}(f_i, f'_i : U_i \to V_{\alpha(i)})
\}_{i \in I,\ \alpha(i) = \alpha'(i)}
$$
Les égalisateurs existent si $\mathcal{C}$ admet des égalisateurs.

\medskip\noindent
Démonstration de (6). Si $X$ est un objet final de $\mathcal{C}$, alors
$\{X\}$ est un objet final de $\text{SR}(\mathcal{C})$.

\medskip\noindent
Démonstration des énoncés concernant $\text{SR}(\mathcal{C}, X)$.
Ils résultent des résultats précédents appliqués à la catégorie
$\mathcal{C}/X$, compte tenu de
$\text{SR}(\mathcal{C}/X) = \text{SR}(\mathcal{C}, X)$ et de
$\textit{PSh}(\mathcal{C}/X) = \textit{PSh}(\mathcal{C})/h_X$
(Sites, lemme \ref{sites-lemma-essential-image-j-shriek}, appliqué
à $\mathcal{C}$ munie de la topologie chaotique). Toutefois,
nous pouvons aussi raisonner directement comme suit.
Il est clair que le coproduit de
$\{U_i \to X\}_{i \in I}$ et $\{V_j \to X\}_{j \in J}$
est $\{U_i \to X\}_{i \in I} \amalg \{V_j \to X\}_{j \in J}$,
et il en va de même pour les coproduits de
familles de familles de morphismes de but $X$.
L'objet $\{X \to X\}$ est un objet
final de $\text{SR}(\mathcal{C}, X)$.
Supposons donnés un morphisme
$(\alpha, f_i) : \{U_i \to X\}_{i \in I} \to \{V_j \to X\}_{j \in J}$
et un morphisme
$(\beta, g_k) : \{W_k \to X\}_{k \in K} \to \{V_j \to X\}_{j \in J}$.
Le produit fibré de ces morphismes est donné par
$$
\{ U_i \times_{f_i, V_j, g_k} W_k \to X \}_{(i, j, k) \in I \times J \times K
\text{ tels que } j = \alpha(i) = \beta(k)}
$$
Les produits fibrés existent par hypothèse, puisque
$\mathcal{C}$ admet des produits fibrés.
Ainsi, $\text{SR}(\mathcal{C}, X)$ admet des limites finies,
voir Catégories, lemme \ref{categories-lemma-finite-limits-exist}.
Nous omettons dans ce cas la vérification des énoncés sur le foncteur $F$.
\end{proof}





\section{Hyperrecouvrements}
\label{section-hypercoverings}

\noindent
Si nous supposons que notre catégorie est un site, nous pouvons poser la
définition suivante.

\begin{definition}
\label{definition-covering-SR}
Soit $\mathcal{C}$ un site. Soit
$f = (\alpha, f_i) : \{U_i\}_{i \in I} \to \{V_j\}_{j \in J}$
un morphisme de la catégorie $\text{SR}(\mathcal{C})$.
Nous disons que $f$ est un {\it recouvrement} si, pour tout $j \in J$, la
famille de morphismes $\{U_i \to V_j\}_{i \in I, \alpha(i) = j}$
est un recouvrement du site $\mathcal{C}$.
Soit $X$ un objet de $\mathcal{C}$.
Un morphisme $K \to L$ dans $\text{SR}(\mathcal{C}, X)$ est
un {\it recouvrement} si son image dans $\text{SR}(\mathcal{C})$ est
un recouvrement.
\end{definition}

\begin{lemma}
\label{lemma-covering-permanence}
Soit $\mathcal{C}$ un site.
\begin{enumerate}
\item Un composé de recouvrements dans $\text{SR}(\mathcal{C})$
est un recouvrement.
\item Si $K \to L$ est un recouvrement dans $\text{SR}(\mathcal{C})$
et si $L' \to L$ est un morphisme, alors $L' \times_L K$ existe
et $L' \times_L K \to L'$ est un recouvrement.
\item Si $\mathcal{C}$ admet les produits de deux objets et si
$A \to B$ et $K \to L$ sont des recouvrements dans $\text{SR}(\mathcal{C})$,
alors $A \times K \to B \times L$ est un recouvrement.
\end{enumerate}
Soit $X \in \Ob(\mathcal{C})$. Alors (1) et (2) valent pour
$\text{SR}(\mathcal{C}, X)$, et (3) vaut si $\mathcal{C}$
admet des produits fibrés.
\end{lemma}

\begin{proof}
La partie (1) découle immédiatement des axiomes d'un site.
La partie (2) résulte de la construction des produits fibrés
dans $\text{SR}(\mathcal{C})$ donnée dans la démonstration du
lemme \ref{lemma-coprod-prod-SR}
et de l'exigence que les morphismes d'un recouvrement
de $\mathcal{C}$ soient représentables.
La partie (3) résulte de la décomposition de $A \times K \to B \times L$
en le composé $A \times K \to B \times K \to B \times L$
et donc en un composé de changements de base de recouvrements.
Le dernier énoncé découle de l'égalité $\text{SR}(\mathcal{C}, X) =
\text{SR}(\mathcal{C}/X)$.
\end{proof}

\noindent
D'après le lemme \ref{lemma-coprod-prod-SR} et
Simplicial, lemme \ref{simplicial-lemma-existence-cosk},
le cosquelette d'un objet simplicial tronqué de
$\text{SR}(\mathcal{C}, X)$ existe si $\mathcal{C}$ admet des produits fibrés.
La définition suivante a donc un sens.

\begin{definition}
\label{definition-hypercovering}
Soit $\mathcal{C}$ un site. Supposons que $\mathcal{C}$ admette des produits fibrés.
Soit $X \in \Ob(\mathcal{C})$ un objet de $\mathcal{C}$.
Un {\it hyperrecouvrement de $X$} est un objet simplicial
$K$ de $\text{SR}(\mathcal{C}, X)$ tel que
\begin{enumerate}
\item l'objet $K_0$ soit un recouvrement de $X$ pour le site $\mathcal{C}$ ;
\item pour tout $n \geq 0$, le morphisme canonique
$$
K_{n + 1} \longrightarrow (\text{cosk}_n \text{sk}_n K)_{n + 1}
$$
soit un recouvrement au sens défini ci-dessus.
\end{enumerate}
\end{definition}

\noindent
La condition (1) a un sens puisque tout objet de
$\text{SR}(\mathcal{C}, X)$ est, après tout, une famille
de morphismes de but $X$. On pourrait aussi
la formuler en disant que le morphisme de $K_0$ vers
l'objet final de $\text{SR}(\mathcal{C}, X)$
est un recouvrement.

\begin{example}[Hyperrecouvrements de {\v C}ech]
\label{example-cech}
Soit $\mathcal{C}$ un site admettant des produits fibrés.
Soit $\{U_i \to X\}_{i \in I}$ un recouvrement de $\mathcal{C}$.
Posons $K_0 = \{U_i \to X\}_{i \in I}$.
Alors $K_0$ est un objet simplicial $0$-tronqué de
$\text{SR}(\mathcal{C}, X)$. Nous pouvons donc former
$$
K = \text{cosk}_0 K_0.
$$
Il est clair que $K$ satisfait la condition (1) de la définition \ref{definition-hypercovering}.
Puisque tous les morphismes $K_{n + 1} \to (\text{cosk}_n \text{sk}_n K)_{n + 1}$
sont des isomorphismes d'après
Simplicial, lemme \ref{simplicial-lemma-cosk-up},
il satisfait aussi la condition (2). Notons que
les termes $K_n$ sont les termes usuels
$$
K_n = \{
U_{i_0} \times_X U_{i_1} \times_X \ldots \times_X U_{i_n} \to X
\}_{(i_0, i_1, \ldots, i_n) \in I^{n + 1}}
$$
Un hyperrecouvrement de $X$ de cette forme est appelé un
{\it hyperrecouvrement de {\v C}ech} de $X$.
\end{example}

\begin{example}[Hyperrecouvrement par un objet simplicial du site]
\label{example-hypercovering-in-C}
Soit $\mathcal{C}$ un site admettant des produits fibrés. Soit
$X \in \Ob(\mathcal{C})$. Soit $U$ un objet simplicial de $\mathcal{C}$.
Comme d'habitude, nous notons $U_n = U([n])$. Supposons enfin donnée une augmentation
$$
a : U \to X
$$
Dans cette situation, nous pouvons considérer l'objet simplicial $K$
de $\text{SR}(\mathcal{C}, X)$ dont les termes sont $K_n = \{U_n \to X\}$.
Alors $K$ est un hyperrecouvrement de $X$ au sens de la
définition \ref{definition-hypercovering}
si et seulement si les trois
conditions suivantes\footnote{Comme $\mathcal{C}$ admet des produits fibrés, la
catégorie $\mathcal{C}/X$ admet toutes les limites finies.
Les cosquelettes requis existent donc d'après
Simplicial, lemme \ref{simplicial-lemma-existence-cosk}.} sont satisfaites :
\begin{enumerate}
\item $\{U_0 \to X\}$ est un recouvrement de $\mathcal{C}$ ;
\item $\{U_1 \to U_0 \times_X U_0\}$ est un recouvrement de $\mathcal{C}$ ;
\item $\{U_{n + 1} \to (\text{cosk}_n\text{sk}_n U)_{n + 1}\}$
est un recouvrement de $\mathcal{C}$ pour $n \geq 1$.
\end{enumerate}
Nous omettons la vérification immédiate.
\end{example}

\begin{example}[Hyperrecouvrement de {\v C}ech associé à un recouvrement]
\label{example-cech-cover}
Soit $\mathcal{C}$ un site admettant des produits fibrés. Soit $U \to X$ un
morphisme de $\mathcal{C}$ tel que $\{U \to X\}$ soit un recouvrement de
$\mathcal{C}$\footnote{Un morphisme de $\mathcal{C}$ ayant cette propriété
est parfois appelé un « recouvrement ».}. Considérons l'objet simplicial $K$ de
$\text{SR}(\mathcal{C}, X)$ dont les termes sont
$$
K_n = \{U \times_X U \times_X \ldots \times_X U \to X\}
\quad (n + 1 \text{ facteurs})
$$
Alors $K$ est un hyperrecouvrement de $X$. Cet exemple est un cas particulier à la fois de
l'exemple \ref{example-cech} et de
l'exemple \ref{example-hypercovering-in-C}.
\end{example}

\begin{lemma}
\label{lemma-hypercoverings-set}
Soit $\mathcal{C}$ un site admettant des produits fibrés.
Soit $X \in \Ob(\mathcal{C})$ un objet de $\mathcal{C}$.
La collection de tous les hyperrecouvrements de $X$ forme un ensemble.
\end{lemma}

\begin{proof}
Puisque $\mathcal{C}$ est un site, l'ensemble de tous les recouvrements de
$X$ est un ensemble. Nous voyons donc que la collection
des $K_0$ possibles forme un ensemble. Supposons avoir montré que
la collection de tous les $K_0, \ldots, K_n$ possibles forme
un ensemble. Il suffit alors de montrer que, étant donnés
$K_0, \ldots, K_n$, la collection de tous les
$K_{n + 1}$ possibles forme un ensemble. Cela est clair puisque
nous devons choisir $K_{n + 1}$ parmi tous les recouvrements possibles
de $(\text{cosk}_n \text{sk}_n K)_{n + 1}$.
\end{proof}

\begin{remark}
\label{remark-hypercoverings-really-set}
Le lemme ne dit pas seulement qu'il existe un système cofinal
de choix d'hyperrecouvrements qui soit un ensemble,
mais bien que les hyperrecouvrements eux-mêmes forment un ensemble.
\end{remark}

\noindent
La catégorie des préfaisceaux sur $\mathcal{C}$ admet des
(co)limites finies. Les foncteurs $\text{cosk}_n$
existent donc pour les préfaisceaux d'ensembles.

\begin{lemma}
\label{lemma-hypercovering-F}
Soit $\mathcal{C}$ un site admettant des produits fibrés.
Soit $X \in \Ob(\mathcal{C})$ un objet de $\mathcal{C}$.
Soit $K$ un hyperrecouvrement de $X$.
Considérons l'objet simplicial $F(K)$ de $\textit{PSh}(\mathcal{C})$,
muni de son augmentation vers le préfaisceau simplicial constant $h_X$.
\begin{enumerate}
\item Le morphisme de préfaisceaux $F(K)_0 \to h_X$ devient
une surjection après faisceautisation.
\item Le morphisme
$$
(d^1_0, d^1_1) :
F(K)_1
\longrightarrow
F(K)_0 \times_{h_X} F(K)_0
$$
devient une surjection après faisceautisation.
\item Pour tout $n \geq 1$, le morphisme
$$
F(K)_{n + 1} \longrightarrow (\text{cosk}_n \text{sk}_n F(K))_{n + 1}
$$
devient une surjection après faisceautisation.
\end{enumerate}
\end{lemma}

\begin{proof}
Nous utiliserons le fait que, si
$\{U_i \to U\}_{i \in I}$ est un recouvrement du site
$\mathcal{C}$, alors le morphisme
$$
\amalg_{i \in I} h_{U_i} \to h_U
$$
devient surjectif après faisceautisation, voir
Sites, lemme \ref{sites-lemma-covering-surjective-after-sheafification}.
La première assertion en résulte immédiatement.

\medskip\noindent
Pour la deuxième assertion, notons que, d'après
Simplicial, exemple \ref{simplicial-example-cosk0},
l'objet simplicial $\text{cosk}_0 \text{sk}_0 K$
a pour termes $K_0 \times \ldots \times K_0$. Ainsi,
d'après la définition d'un hyperrecouvrement, nous
voyons que $(d^1_0, d^1_1) : K_1 \to K_0 \times K_0$ est un
recouvrement. Par conséquent, (2) résulte de l'assertion précédente
et du fait que $F$ transforme les produits en produits
fibrés au-dessus de $h_X$.

\medskip\noindent
Pour la troisième, nous affirmons que
$\text{cosk}_n \text{sk}_n F(K) =
F(\text{cosk}_n \text{sk}_n K)$ pour $n \geq 1$.
Pour le démontrer, notons provisoirement $F'$ le foncteur
$\text{SR}(\mathcal{C}, X) \to \textit{PSh}(\mathcal{C})/h_X$.
Par le lemme \ref{lemma-coprod-prod-SR}, le foncteur
$F'$ commute aux limites finies.
D'après notre description du foncteur $\text{cosk}_n$ dans
Simplicial, section \ref{simplicial-section-skeleton},
nous voyons que $\text{cosk}_n \text{sk}_n F'(K) =
F'(\text{cosk}_n \text{sk}_n K)$.
Rappelons que la catégorie utilisée dans la description de
$(\text{cosk}_n U)_m$ dans
Simplicial, lemme \ref{simplicial-lemma-existence-cosk},
est la catégorie $(\Delta/[m])^{opp}_{\leq n}$. Un exercice
amusant consiste à montrer que $(\Delta/[m])_{\leq n}$ est
une catégorie connexe (voir
Catégories, définition \ref{categories-definition-category-connected})
dès que $n \geq 1$. Par conséquent,
Catégories, lemme \ref{categories-lemma-connected-limit-over-X},
montre que $\text{cosk}_n \text{sk}_n F'(K) =
\text{cosk}_n \text{sk}_n F(K)$. D'où l'assertion.
La propriété (3) en résulte, car nous voyons maintenant que
le morphisme de (3) s'obtient en appliquant le
foncteur $F$ à un recouvrement comme dans la définition \ref{definition-covering-SR},
et le résultat découle du premier fait mentionné
dans cette démonstration.
\end{proof}



\section{Acyclicité}
\label{section-acyclicity}

\noindent
Soit $\mathcal{C}$ un site.
Pour un préfaisceau d'ensembles $\mathcal{F}$, nous notons $\mathbf{Z}_\mathcal{F}$
le préfaisceau de groupes abéliens défini par la règle
$$
\mathbf{Z}_\mathcal{F}(U) = \text{groupe abélien libre sur }\mathcal{F}(U).
$$
Nous l'appellerons parfois le {\it préfaisceau abélien libre sur $\mathcal{F}$}.
Bien entendu, la construction $\mathcal{F} \mapsto \mathbf{Z}_\mathcal{F}$
est un foncteur, adjoint à gauche du foncteur d'oubli
$\textit{PAb}(\mathcal{C}) \to \textit{PSh}(\mathcal{C})$.
Bien entendu, le faisceau associé $\mathbf{Z}_\mathcal{F}^\#$ est
un faisceau de groupes abéliens, et le foncteur
$\mathcal{F} \mapsto \mathbf{Z}_\mathcal{F}^\#$ est lui aussi
un adjoint à gauche. Nous appelons parfois $\mathbf{Z}_\mathcal{F}^\#$
le {\it faisceau abélien libre sur $\mathcal{F}$}.

\medskip\noindent
Pour un objet $X$ du site $\mathcal{C}$, nous notons
$\mathbf{Z}_X$ le préfaisceau abélien libre sur $h_X$, et
nous notons $\mathbf{Z}_X^\#$ son faisceau associé.

\begin{definition}
\label{definition-homology}
Soit $\mathcal{C}$ un site.
Soit $K$ un objet simplicial de $\textit{PSh}(\mathcal{C})$.
D'après ce qui précède, nous obtenons un objet simplicial $\mathbf{Z}_K^\#$ de
$\textit{Ab}(\mathcal{C})$. Nous pouvons prendre son
complexe associé de faisceaux abéliens $s(\mathbf{Z}_K^\#)$, voir
Simplicial, section \ref{simplicial-section-complexes}.
L'{\it homologie de $K$} est l'homologie du
complexe de faisceaux abéliens $s(\mathbf{Z}_K^\#)$.
\end{definition}

\noindent
Autrement dit, la {\it $i$-ième homologie $H_i(K)$ de $K$}
est le faisceau de groupes abéliens $H_i(K) = H_i(s(\mathbf{Z}_K^\#))$.
Dans cette section, nous étudions l'homologie lorsque $K$
est un hyperrecouvrement d'un objet $X$ de $\mathcal{C}$.

\begin{lemma}
\label{lemma-compare-cosk0}
Soit $\mathcal{C}$ un site.
Soit $\mathcal{F} \to \mathcal{G}$ un morphisme
de préfaisceaux d'ensembles. Notons $K$ l'objet
simplicial de $\textit{PSh}(\mathcal{C})$ dont le $n$-ième
terme est le produit fibré à $(n + 1)$ facteurs de $\mathcal{F}$
au-dessus de $\mathcal{G}$, voir
Simplicial, exemple \ref{simplicial-example-fibre-products-simplicial-object}.
Alors, si $\mathcal{F} \to \mathcal{G}$ est surjectif après
faisceautisation, nous avons
$$
H_i(K) =
\left\{
\begin{matrix}
0 & \text{si} & i > 0\\
\mathbf{Z}_\mathcal{G}^\# & \text{si} & i = 0
\end{matrix}
\right.
$$
L'isomorphisme en degré $0$ est donné par le
morphisme $H_0(K) \to \mathbf{Z}_\mathcal{G}^\#$
provenant de l'application $(\mathbf{Z}_K^\#)_0 =
\mathbf{Z}_\mathcal{F}^\# \to \mathbf{Z}_\mathcal{G}^\#$.
\end{lemma}

\begin{proof}
Soit $\mathcal{G}' \subset \mathcal{G}$ l'image du
morphisme $\mathcal{F} \to \mathcal{G}$.
Soit $U \in \Ob(\mathcal{C})$. Posons
$A = \mathcal{F}(U)$ et $B = \mathcal{G}'(U)$.
Alors l'ensemble simplicial $K(U)$ est égal à l'ensemble
simplicial dont les $n$-simplexes sont donnés par
$$
A \times_B A \times_B \ldots \times_B A\ (n + 1 \text{ facteurs)}.
$$
D'après Simplicial, lemme \ref{simplicial-lemma-cosk-minus-one-equivalence},
le morphisme $K(U) \to B$ est une fibration de Kan triviale.
C'est donc une équivalence d'homotopie
(Simplicial, lemme \ref{simplicial-lemma-trivial-kan-homotopy}).
En appliquant à celle-ci le foncteur « groupe abélien libre sur »,
nous en déduisons que
$$
\mathbf{Z}_K(U) \longrightarrow \mathbf{Z}_B
$$
est une équivalence d'homotopie. Notons que $s(\mathbf{Z}_B)$ est
le complexe
$$
\ldots \to
\bigoplus\nolimits_{b \in B}\mathbf{Z} \xrightarrow{0}
\bigoplus\nolimits_{b \in B}\mathbf{Z} \xrightarrow{1}
\bigoplus\nolimits_{b \in B}\mathbf{Z} \xrightarrow{0}
\bigoplus\nolimits_{b \in B}\mathbf{Z} \to 0
$$
voir Simplicial, lemme \ref{simplicial-lemma-homology-eilenberg-maclane}.
Nous voyons ainsi que
$H_i(s(\mathbf{Z}_K(U))) = 0$ pour $i > 0$, et
$H_0(s(\mathbf{Z}_K(U))) = \bigoplus_{b \in B}\mathbf{Z}
= \bigoplus_{s \in \mathcal{G}'(U)} \mathbf{Z}$.
Ces identifications sont compatibles aux applications de
restriction.

\medskip\noindent
Nous concluons que $H_i(s(\mathbf{Z}_K)) = 0$ pour $i > 0$ et
$H_0(s(\mathbf{Z}_K)) = \mathbf{Z}_{\mathcal{G}'}$, où nous
calculons ici les groupes d'homologie dans $\textit{PAb}(\mathcal{C})$. Comme
la faisceautisation est un foncteur exact, nous en déduisons le résultat
du lemme. En effet, l'exactitude implique
que $H_0(s(\mathbf{Z}_K))^\# = H_0(s(\mathbf{Z}_K^\#))$,
et de même pour les autres indices.
\end{proof}

\begin{lemma}
\label{lemma-acyclicity}
Soit $\mathcal{C}$ un site.
Soit $f : L \to K$ un morphisme d'objets
simpliciaux de $\textit{PSh}(\mathcal{C})$.
Soit $n \geq 0$ un entier.
Supposons que
\begin{enumerate}
\item pour $i < n$, le morphisme $L_i \to K_i$ soit un isomorphisme ;
\item le morphisme $L_n \to K_n$ soit surjectif après faisceautisation ;
\item l'application canonique $L \to \text{cosk}_n \text{sk}_n L$ soit un isomorphisme ;
\item l'application canonique $K \to \text{cosk}_n \text{sk}_n K$ soit un isomorphisme.
\end{enumerate}
Alors $H_i(f) : H_i(L) \to H_i(K)$ est un isomorphisme.
\end{lemma}

\begin{proof}
Cette démonstration est exactement la même que celle du
lemme \ref{lemma-compare-cosk0} ci-dessus. Plus précisément,
soit d'abord $K_n' \subset K_n$ le sous-préfaisceau
image de l'application $L_n \to K_n$. L'hypothèse
(2) signifie que le faisceau associé à $K_n'$ est égal au
faisceau associé à $K_n$. En outre, puisque $L_i = K_i$
pour tout $i < n$, nous obtenons un préfaisceau simplicial
$n$-tronqué $U$ en posant
$U_0 = L_0 = K_0, \ldots, U_{n - 1} = L_{n - 1} = K_{n - 1}, U_n = K'_n$.
Notons $K' = \text{cosk}_n U$, qui est un préfaisceau simplicial.
Comme nous pouvons construire $K'_m$ comme une limite finie et
que la faisceautisation est exacte, nous voyons que
$(K'_m)^\# = K_m$. Autrement dit, $(K')^\# = K^\#$.
Nous concluons, de nouveau par exactitude de la faisceautisation,
que $H_i(K) = H_i(K')$. Il suffit donc de démontrer le lemme
pour le morphisme $L \to K'$ ; autrement dit, nous pouvons
supposer que $L_n \to K_n$ est un morphisme surjectif
de {\it préfaisceaux} !

\medskip\noindent
Dans ce cas, pour tout objet $U$ de $\mathcal{C}$, nous
voyons que le morphisme d'ensembles simpliciaux
$$
L(U) \longrightarrow K(U)
$$
satisfait toutes les hypothèses de
Simplicial, lemme \ref{simplicial-lemma-section}.
C'est donc une fibration de Kan triviale. En particulier, c'est
une équivalence d'homotopie
(Simplicial, lemme \ref{simplicial-lemma-trivial-kan-homotopy}).
Ainsi,
$$
\mathbf{Z}_L(U) \longrightarrow \mathbf{Z}_K(U)
$$
est aussi une équivalence d'homotopie. Cela vaut pour tout $U$.
Le résultat en découle.
\end{proof}

\begin{lemma}
\label{lemma-acyclic-hypercover-sheaves}
Soit $\mathcal{C}$ un site.
Soit $K$ un préfaisceau simplicial.
Soit $\mathcal{G}$ un préfaisceau.
Soit $K \to \mathcal{G}$ une augmentation de $K$
vers $\mathcal{G}$. Supposons que
\begin{enumerate}
\item le morphisme de préfaisceaux $K_0 \to \mathcal{G}$ devienne
une surjection après faisceautisation ;
\item le morphisme
$$
(d^1_0, d^1_1) :
K_1
\longrightarrow
K_0 \times_\mathcal{G} K_0
$$
devienne une surjection après faisceautisation ;
\item pour tout $n \geq 1$, le morphisme
$$
K_{n + 1} \longrightarrow (\text{cosk}_n \text{sk}_n K)_{n + 1}
$$
devienne une surjection après faisceautisation.
\end{enumerate}
Alors $H_i(K) = 0$ pour $i > 0$ et
$H_0(K) = \mathbf{Z}_\mathcal{G}^\#$.
\end{lemma}

\begin{proof}
Notons $K^n = \text{cosk}_n \text{sk}_n K$ pour $n \geq 1$.
Définissons $K^0$ comme l'objet simplicial dont les termes
$(K^0)_n$ sont égaux au produit fibré à $(n + 1)$ facteurs
$K_0 \times_\mathcal{G} \ldots \times_\mathcal{G} K_0$,
voir Simplicial,
exemple \ref{simplicial-example-fibre-products-simplicial-object}.
Nous avons des morphismes
$$
K \longrightarrow \ldots \to K^n \to K^{n - 1} \to \ldots \to K^1 \to K^0.
$$
Les morphismes $K \to K^i$, $K^j \to K^i$ pour $j \geq i \geq 1$ proviennent
des propriétés universelles des foncteurs $\text{cosk}_n$.
Le morphisme $K^1 \to K^0$ est le morphisme canonique
de
Simplicial, remarque \ref{simplicial-remark-augmentation}.
Rappelons aussi que $K^0 \to \text{cosk}_1 \text{sk}_1 K^0$
est un isomorphisme, voir
Simplicial, lemme \ref{simplicial-lemma-cosk-minus-one}.

\medskip\noindent
D'après le lemme \ref{lemma-compare-cosk0}, nous voyons que
$H_i(K^0) = 0$ pour $i > 0$ et $H_0(K^0) = \mathbf{Z}_\mathcal{G}^\#$.

\medskip\noindent
Fixons $n \geq 1$. Considérons le morphisme $K^n \to K^{n - 1}$.
C'est un isomorphisme sur les termes de degré $< n$.
Notons que $K^n \to \text{cosk}_n \text{sk}_n K^n$ et
$K^{n - 1} \to \text{cosk}_n \text{sk}_n K^{n - 1}$
sont des isomorphismes. Notons que $(K^n)_n = K_n$ et
que $(K^{n - 1})_n = (\text{cosk}_{n - 1} \text{sk}_{n - 1} K)_n$.
Par hypothèse, $(K^n)_n \to (K^{n - 1})_n$
est donc un morphisme de préfaisceaux qui devient surjectif après
faisceautisation. Par le lemme \ref{lemma-acyclicity}, nous concluons que
$H_i(K^n) = H_i(K^{n - 1})$.
Combiné à ce qui précède, cela démontre le lemme.
\end{proof}

\begin{lemma}
\label{lemma-hypercovering-acyclic}
Soit $\mathcal{C}$ un site admettant des produits fibrés.
Soit $X$ un objet de $\mathcal{C}$.
Soit $K$ un hyperrecouvrement de $X$.
L'homologie du préfaisceau simplicial $F(K)$ est
$0$ dans les degrés $> 0$ et égale à $\mathbf{Z}_X^\#$
en degré $0$.
\end{lemma}

\begin{proof}
Combiner les lemmes \ref{lemma-acyclic-hypercover-sheaves}
et \ref{lemma-hypercovering-F}.
\end{proof}












\section{Cohomologie de {\v C}ech et hyperrecouvrements}
\label{section-hyper-cech}

\noindent
Soit $\mathcal{C}$ un site. Considérons un préfaisceau de
groupes abéliens $\mathcal{F}$ sur le site $\mathcal{C}$.
Il définit un foncteur
\begin{eqnarray*}
\mathcal{F} : \text{SR}(\mathcal{C})^{opp}
& \longrightarrow &
\textit{Ab} \\
\{U_i\}_{i \in I} &
\longmapsto &
\prod\nolimits_{i \in I} \mathcal{F}(U_i)
\end{eqnarray*}
Ainsi, un objet simplicial $K$ de $\text{SR}(\mathcal{C})$
est transformé en un objet cosimplicial $\mathcal{F}(K)$ de $\textit{Ab}$.
Le complexe de cochaînes $s(\mathcal{F}(K))$ associé à $\mathcal{F}(K)$
(Simplicial, section
\ref{simplicial-section-dold-kan-cosimplicial})
est appelé le complexe de {\v C}ech de $\mathcal{F}$ relativement
à l'objet simplicial $K$. Nous posons
$$
\check{H}^i(K, \mathcal{F})
=
H^i(s(\mathcal{F}(K))).
$$
et nous l'appelons le $i$-ième groupe de cohomologie de {\v C}ech
de $\mathcal{F}$ relativement à $K$.
Dans cette section, nous démontrons des analogues de certains résultats sur
la cohomologie de {\v C}ech des recouvrements ouverts établis dans
Cohomologie, sections \ref{cohomology-section-cech},
\ref{cohomology-section-cech-functor} et
\ref{cohomology-section-cech-cohomology-cohomology}.

\begin{lemma}
\label{lemma-h0-cech}
Soit $\mathcal{C}$ un site admettant des produits fibrés.
Soit $X$ un objet de $\mathcal{C}$.
Soit $K$ un hyperrecouvrement de $X$.
Soit $\mathcal{F}$ un faisceau de groupes abéliens sur $\mathcal{C}$.
Alors $\check{H}^0(K, \mathcal{F}) = \mathcal{F}(X)$.
\end{lemma}

\begin{proof}
Nous avons
$$
\check{H}^0(K, \mathcal{F})
=
\Ker(\mathcal{F}(K_0) \longrightarrow \mathcal{F}(K_1))
$$
Écrivons $K_0 = \{U_i \to X\}$. C'est un recouvrement du site
$\mathcal{C}$. De plus, $K_1 \to K_0 \times K_0$
est un recouvrement dans $\text{SR}(\mathcal{C}, X)$. Nous pouvons donc
écrire $K_1 = \amalg_{i_0, i_1 \in I} \{V_{i_0i_1j} \to X\}$
de telle sorte que le morphisme $K_1 \to K_0 \times K_0$ soit donné
par les recouvrements $\{V_{i_0i_1j} \to U_{i_0} \times_X U_{i_1}\}$
du site $\mathcal{C}$. On obtient donc l’identification suivante :
$$
\check{H}^0(K, \mathcal{F})
=
\Ker(
\prod\nolimits_i \mathcal{F}(U_i)
\longrightarrow
\prod\nolimits_{i_0i_1 j} \mathcal{F}(V_{i_0i_1j})
)
$$
au moyen de l'application évidente. La propriété de faisceau de $\mathcal{F}$
implique que $\check{H}^0(K, \mathcal{F}) = H^0(X, \mathcal{F})$.
\end{proof}

\noindent
En fait, cette propriété caractérise bien entendu les faisceaux abéliens parmi tous
les préfaisceaux abéliens sur $\mathcal{C}$.
L'analogue de Cohomologie, lemme \ref{lemma-injective-trivial-cech},
est dans ce cas le suivant.

\begin{lemma}
\label{lemma-injective-trivial-cech}
Soit $\mathcal{C}$ un site admettant des produits fibrés.
Soit $X$ un objet de $\mathcal{C}$.
Soit $K$ un hyperrecouvrement de $X$.
Soit $\mathcal{I}$ un faisceau injectif de groupes abéliens sur $\mathcal{C}$.
Alors
$$
\check{H}^p(K, \mathcal{I}) =
\left\{
\begin{matrix}
\mathcal{I}(X) & \text{si} & p = 0 \\
0 & \text{si} & p > 0
\end{matrix}
\right.
$$
\end{lemma}

\begin{proof}
Observons que, pour tout objet $Z = \{U_i \to X\}$ de
$\text{SR}(\mathcal{C}, X)$ et tout faisceau abélien
$\mathcal{F}$ sur $\mathcal{C}$, nous avons
\begin{eqnarray*}
\mathcal{F}(Z)
& = &
\prod \mathcal{F}(U_i) \\
& = &
\prod \Mor_{\textit{PSh}(\mathcal{C})}(h_{U_i}, \mathcal{F})\\
& = &
\Mor_{\textit{PSh}(\mathcal{C})}(F(Z), \mathcal{F})\\
& = &
\Mor_{\textit{PAb}(\mathcal{C})}(\mathbf{Z}_{F(Z)}, \mathcal{F}) \\
& = &
\Mor_{\textit{Ab}(\mathcal{C})}(\mathbf{Z}_{F(Z)}^\#, \mathcal{F})
\end{eqnarray*}
Ainsi, nous voyons que, pour tout objet simplicial $K$ de
$\text{SR}(\mathcal{C}, X)$, nous avons
\begin{equation}
\label{equation-identify-cech}
s(\mathcal{F}(K))
=
\Hom_{\textit{Ab}(\mathcal{C})}(s(\mathbf{Z}_{F(K)}^\#), \mathcal{F})
\end{equation}
voir la définition \ref{definition-homology} pour les notations.
Le complexe de faisceaux $s(\mathbf{Z}_{F(K)}^\#)$ est quasi-isomorphe
à $\mathbf{Z}_X^\#$ si $K$ est un hyperrecouvrement, voir le
lemme \ref{lemma-hypercovering-acyclic}. Nous concluons
que, si $\mathcal{I}$ est un faisceau abélien injectif et
$K$ un hyperrecouvrement, alors le complexe $s(\mathcal{I}(K))$
est acyclique, sauf peut-être en degré $0$.
Autrement dit, nous avons
$$
\check{H}^i(K, \mathcal{I}) = 0
$$
pour $i > 0$. Combiné au lemme \ref{lemma-h0-cech}, cela démontre le lemme.
\end{proof}

\noindent
Passons maintenant à l'analogue de Cohomologie sur les sites, lemme
\ref{sites-cohomology-lemma-cech-spectral-sequence}.
Soit $\mathcal{C}$ un site.
Soit $\mathcal{F}$ un faisceau de groupes abéliens sur $\mathcal{C}$.
Rappelons que $\underline{H}^i(\mathcal{F})$ désigne le préfaisceau
de groupes abéliens sur $\mathcal{C}$ défini par la
règle $\underline{H}^i(\mathcal{F}) : U \longmapsto H^i(U, \mathcal{F})$.
Nous l'étendons à $\text{SR}(\mathcal{C})$ comme dans l'introduction
de cette section.

\begin{lemma}
\label{lemma-cech-spectral-sequence}
Soit $\mathcal{C}$ un site admettant des produits fibrés.
Soit $X$ un objet de $\mathcal{C}$.
Soit $K$ un hyperrecouvrement de $X$.
Soit $\mathcal{F}$ un faisceau de groupes abéliens sur $\mathcal{C}$.
Il existe une application
$$
s(\mathcal{F}(K))
\longrightarrow
R\Gamma(X, \mathcal{F})
$$
dans $D^{+}(\textit{Ab})$, fonctorielle en $\mathcal{F}$, qui induit des
transformations naturelles
$$
\check{H}^i(K, -) \longrightarrow H^i(X, -)
$$
de foncteurs $\textit{Ab}(\mathcal{C}) \to \textit{Ab}$. En outre,
il existe une suite spectrale $(E_r, d_r)_{r \geq 0}$ avec
$$
E_2^{p, q} = \check{H}^p(K, \underline{H}^q(\mathcal{F}))
$$
qui converge vers $H^{p + q}(X, \mathcal{F})$.
Cette suite spectrale est fonctorielle en $\mathcal{F}$ et
en l'hyperrecouvrement $K$.
\end{lemma}

\begin{proof}
Nous pourrions démontrer ceci par la même méthode que celle employée dans le lemme
correspondant du chapitre sur la cohomologie. Donnons plutôt un argument
par complexe double.

\medskip\noindent
Choisissons une résolution injective $\mathcal{F} \to \mathcal{I}^\bullet$
dans la catégorie des faisceaux abéliens sur $\mathcal{C}$. Considérons le
complexe double $A^{\bullet, \bullet}$ de termes
$$
A^{p, q} = \mathcal{I}^q(K_p)
$$
où la différentielle $d_1^{p, q} : A^{p, q} \to A^{p + 1, q}$ provient
de la différentielle du complexe $s(\mathcal{I}^q(K))$
associé au groupe abélien cosimplicial $\mathcal{I}^p(K)$,
et où la différentielle $d_2^{p, q} : A^{p, q} \to A^{p, q + 1}$ provient
de la différentielle $\mathcal{I}^q \to \mathcal{I}^{q + 1}$.
Notons $\text{Tot}(A^{\bullet, \bullet})$ le complexe total associé au
complexe double $A^{\bullet, \bullet}$, voir
Homologie, section \ref{homology-section-double-complexes}.
Nous utiliserons les deux suites
spectrales $({}'E_r, {}'d_r)$ et $({}''E_r, {}''d_r)$
associées à ce complexe double, voir
Homologie, section \ref{homology-section-double-complex}.

\medskip\noindent
D'après le lemme \ref{lemma-injective-trivial-cech},
les complexes $s(\mathcal{I}^q(K))$ sont acycliques en
degrés positifs et ont leur $H^0$ égal à $\mathcal{I}^q(X)$.
Par conséquent, d'après
Homologie, lemme \ref{homology-lemma-double-complex-gives-resolution},
l'application naturelle
$$
\mathcal{I}^\bullet(X) \longrightarrow \text{Tot}(A^{\bullet, \bullet})
$$
est un quasi-isomorphisme de complexes de groupes abéliens. En particulier,
nous concluons que $H^n(\text{Tot}(A^{\bullet, \bullet})) = H^n(X, \mathcal{F})$.

\medskip\noindent
L'application $s(\mathcal{F}(K)) \longrightarrow R\Gamma(X, \mathcal{F})$ du
lemme est le composé de l'application
$s(\mathcal{F}(K)) \to \text{Tot}(A^{\bullet, \bullet})$
et de l'inverse
du quasi-isomorphisme affiché ci-dessus. Cela convient car
$\mathcal{I}^\bullet(X)$ est un représentant de $R\Gamma(X, \mathcal{F})$.

\medskip\noindent
Considérons la suite spectrale $({}'E_r, {}'d_r)_{r \geq 0}$. D'après
Homologie, lemme \ref{homology-lemma-ss-double-complex},
nous voyons que
$$
{}'E_2^{p, q} = H^p_I(H^q_{II}(A^{\bullet, \bullet}))
$$
Autrement dit, nous prenons d'abord la cohomologie relativement à
$d_2$, ce qui donne les groupes
${}'E_1^{p, q} = \underline{H}^q(\mathcal{F})(K_p)$.
Ainsi, nous avons bien (d'après la description de la différentielle
${}'d_1$)
${}'E_2^{p, q} = \check{H}^p(K, \underline{H}^q(\mathcal{F}))$.
D'après ce qui précède et Homologie, lemme \ref{homology-lemma-first-quadrant-ss},
nous voyons que cette suite converge vers $H^n(X, \mathcal{F})$, comme souhaité.

\medskip\noindent
Nous omettons la démonstration des énoncés concernant la fonctorialité des
constructions précédentes par rapport au faisceau abélien $\mathcal{F}$ et à
l'hyperrecouvrement $K$.
\end{proof}









\section{Hyperrecouvrements à la Verdier}
\label{section-hypercoverings-verdier}

\noindent
Le lecteur attentif aura remarqué que, pour obtenir la suite spectrale
de {\v C}ech vers la cohomologie associée à un hyperrecouvrement d'un objet $X$,
nous n'avons besoin que de la
conclusion du lemme \ref{lemma-hypercovering-F}.
La définition suivante a donc un sens.

\begin{definition}
\label{definition-hypercovering-variant}
Soit $\mathcal{C}$ un site. Supposons que $\mathcal{C}$ admette des égalisateurs
et des produits fibrés. Soit $\mathcal{G}$ un préfaisceau d'ensembles.
Un {\it hyperrecouvrement de $\mathcal{G}$} est un objet simplicial
$K$ de $\text{SR}(\mathcal{C})$ muni d'une augmentation
$F(K) \to \mathcal{G}$ telle que
\begin{enumerate}
\item $F(K_0) \to \mathcal{G}$ devienne surjectif
après faisceautisation ;
\item $F(K_1) \to F(K_0) \times_\mathcal{G} F(K_0)$
devienne surjectif après faisceautisation ;
\item $F(K_{n + 1}) \longrightarrow F((\text{cosk}_n \text{sk}_n K)_{n + 1})$
devienne surjectif après faisceautisation pour $n \geq 1$.
\end{enumerate}
Nous disons qu'un objet simplicial $K$ de $\text{SR}(\mathcal{C})$
est un {\it hyperrecouvrement} si $K$ est un hyperrecouvrement de l'objet
final $*$ de $\textit{PSh}(\mathcal{C})$.
\end{definition}

\noindent
L'hypothèse selon laquelle $\mathcal{C}$ admet des produits fibrés et des égalisateurs
garantit que $\text{SR}(\mathcal{C})$ admet des produits fibrés
et des égalisateurs, et que $F$ commute à ceux-ci
(lemme \ref{lemma-coprod-prod-SR}), ce qui suffit
à définir les foncteurs cosquelette utilisés (voir
Simplicial, remarque \ref{simplicial-remark-existence-cosk} et
Catégories, lemme \ref{categories-lemma-fibre-products-equalizers-exist}).
Si $\mathcal{C}$ est quelconque, nous pouvons remplacer la condition (3) par la
condition selon laquelle
$F(K_{n + 1}) \longrightarrow ((\text{cosk}_n \text{sk}_n F(K))_{n + 1})$
devient surjectif après faisceautisation pour $n \geq 1$ ; les
résultats de cette section restent alors valables.

\medskip\noindent
Soit $\mathcal{F}$ un faisceau abélien sur $\mathcal{C}$.
Dans la section précédente, nous avons défini le complexe de {\v C}ech de $\mathcal{F}$
relativement à un objet simplicial $K$ de $\text{SR}(\mathcal{C})$.
Ensuite, étant donné un préfaisceau $\mathcal{G}$, nous posons
$$
H^0(\mathcal{G}, \mathcal{F}) =
\Mor_{\textit{PSh}(\mathcal{C})}(\mathcal{G}, \mathcal{F}) =
\Mor_{\Sh(\mathcal{C})}(\mathcal{G}^\#, \mathcal{F}) =
H^0(\mathcal{G}^\#, \mathcal{F})
$$
avec les notations de
Cohomologie sur les sites, section \ref{sites-cohomology-section-limp}.
C'est un foncteur exact à gauche, et ses foncteurs dérivés supérieurs
(brièvement étudiés dans
Cohomologie sur les sites, section \ref{sites-cohomology-section-limp})
sont notés $H^i(\mathcal{G}, \mathcal{F})$.
Nous montrerons que, étant donné un hyperrecouvrement $K$ de $\mathcal{G}$,
il existe une suite spectrale de {\v C}ech vers la cohomologie qui converge vers la
cohomologie $H^i(\mathcal{G}, \mathcal{F})$.
Notons que, si $\mathcal{G} = *$, alors
$H^i(*, \mathcal{F}) = H^i(\mathcal{C}, \mathcal{F})$ redonne
la cohomologie de $\mathcal{F}$ sur le site $\mathcal{C}$.

\begin{lemma}
\label{lemma-h0-cech-variant}
Soit $\mathcal{C}$ un site admettant des égalisateurs et des produits fibrés.
Soit $\mathcal{G}$ un préfaisceau sur $\mathcal{C}$.
Soit $K$ un hyperrecouvrement de $\mathcal{G}$.
Soit $\mathcal{F}$ un faisceau de groupes abéliens sur $\mathcal{C}$.
Alors $\check{H}^0(K, \mathcal{F}) = H^0(\mathcal{G}, \mathcal{F})$.
\end{lemma}

\begin{proof}
Cela découle de la définition de $H^0(\mathcal{G}, \mathcal{F})$
et du fait que
$$
\xymatrix{
F(K_1) \ar@<1ex>[r] \ar@<-1ex>[r] &
F(K_0) \ar[r] & \mathcal{G}
}
$$
devient un diagramme de coégalisateur après faisceautisation.
\end{proof}

\begin{lemma}
\label{lemma-injective-trivial-cech-variant}
Soit $\mathcal{C}$ un site admettant des égalisateurs et des produits fibrés.
Soit $\mathcal{G}$ un préfaisceau sur $\mathcal{C}$.
Soit $K$ un hyperrecouvrement de $\mathcal{G}$.
Soit $\mathcal{I}$ un faisceau injectif de groupes abéliens sur $\mathcal{C}$.
Alors
$$
\check{H}^p(K, \mathcal{I}) =
\left\{
\begin{matrix}
H^0(\mathcal{G}, \mathcal{I}) & \text{si} & p = 0 \\
0 & \text{si} & p > 0
\end{matrix}
\right.
$$
\end{lemma}

\begin{proof}
D'après (\ref{equation-identify-cech}), nous avons
$$
s(\mathcal{F}(K))
=
\Hom_{\textit{Ab}(\mathcal{C})}(s(\mathbf{Z}_{F(K)}^\#), \mathcal{F})
$$
Le complexe $s(\mathbf{Z}_{F(K)}^\#)$ est quasi-isomorphe
à $\mathbf{Z}_\mathcal{G}^\#$, voir le
lemme \ref{lemma-acyclic-hypercover-sheaves}. Nous concluons
que, si $\mathcal{I}$ est un faisceau abélien injectif, alors
le complexe $s(\mathcal{I}(K))$ est acyclique, sauf peut-être en degré $0$.
Autrement dit, nous avons $\check{H}^i(K, \mathcal{I}) = 0$
pour $i > 0$. Combiné au lemme \ref{lemma-h0-cech-variant},
cela démontre le lemme.
\end{proof}

\begin{lemma}
\label{lemma-cech-spectral-sequence-variant}
Soit $\mathcal{C}$ un site admettant des égalisateurs et des produits fibrés.
Soit $\mathcal{G}$ un préfaisceau sur $\mathcal{C}$.
Soit $K$ un hyperrecouvrement de $\mathcal{G}$.
Soit $\mathcal{F}$ un faisceau de groupes abéliens sur $\mathcal{C}$.
Il existe une application
$$
s(\mathcal{F}(K)) \longrightarrow R\Gamma(\mathcal{G}, \mathcal{F})
$$
dans $D^{+}(\textit{Ab})$, fonctorielle en $\mathcal{F}$, qui induit
une transformation naturelle
$$
\check{H}^i(K, -) \longrightarrow H^i(\mathcal{G}, -)
$$
de foncteurs $\textit{Ab}(\mathcal{C}) \to \textit{Ab}$. En outre,
il existe une suite spectrale $(E_r, d_r)_{r \geq 0}$ avec
$$
E_2^{p, q} = \check{H}^p(K, \underline{H}^q(\mathcal{F}))
$$
qui converge vers $H^{p + q}(\mathcal{G}, \mathcal{F})$.
Cette suite spectrale est fonctorielle en $\mathcal{F}$ et
en l'hyperrecouvrement $K$.
\end{lemma}

\begin{proof}
Choisissons une résolution injective $\mathcal{F} \to \mathcal{I}^\bullet$
dans la catégorie des faisceaux abéliens sur $\mathcal{C}$. Considérons le
complexe double $A^{\bullet, \bullet}$ de termes
$$
A^{p, q} = \mathcal{I}^q(K_p)
$$
où la différentielle $d_1^{p, q} : A^{p, q} \to A^{p + 1, q}$
est celle provenant de la différentielle $\mathcal{I}^p \to \mathcal{I}^{p + 1}$,
et où la différentielle $d_2^{p, q} : A^{p, q} \to A^{p, q + 1}$ est
celle provenant de la différentielle du complexe
$s(\mathcal{I}^p(K))$ associé au groupe abélien cosimplicial
$\mathcal{I}^p(K)$ comme expliqué ci-dessus.
Nous utiliserons les deux suites
spectrales $({}'E_r, {}'d_r)$ et $({}''E_r, {}''d_r)$
associées à ce complexe double, voir
Homologie, section \ref{homology-section-double-complex}.

\medskip\noindent
D'après le lemme \ref{lemma-injective-trivial-cech-variant}, les complexes
$s(\mathcal{I}^p(K))$ sont acycliques en degrés positifs et ont
leur $H^0$ égal à $H^0(\mathcal{G}, \mathcal{I}^p)$. Par conséquent, d'après
Homologie, lemme \ref{homology-lemma-double-complex-gives-resolution}
et sa démonstration, la suite spectrale $({}'E_r, {}'d_r)$ dégénère,
et l'application naturelle
$$
H^0(\mathcal{G}, \mathcal{I}^\bullet) \longrightarrow
\text{Tot}(A^{\bullet, \bullet})
$$
est un quasi-isomorphisme de complexes de groupes abéliens. L'application
$s(\mathcal{F}(K)) \longrightarrow R\Gamma(\mathcal{G}, \mathcal{F})$
du lemme est le composé de l'application naturelle
$s(\mathcal{F}(K)) \to \text{Tot}(A^{\bullet, \bullet})$
et de l'inverse du quasi-isomorphisme affiché ci-dessus.
Cela convient car $H^0(\mathcal{G}, \mathcal{I}^\bullet)$
est un représentant de $R\Gamma(\mathcal{G}, \mathcal{F})$.

\medskip\noindent
Considérons la suite spectrale $({}''E_r, {}''d_r)_{r \geq 0}$. D'après
Homologie, lemme \ref{homology-lemma-ss-double-complex},
nous voyons que
$$
{}''E_2^{p, q} = H^p_{II}(H^q_I(A^{\bullet, \bullet}))
$$
Autrement dit, nous prenons d'abord la cohomologie relativement à
$d_1$, ce qui donne les groupes
${}''E_1^{p, q} = \underline{H}^p(\mathcal{F})(K_q)$.
Ainsi, nous avons bien (d'après la description de la différentielle
${}''d_1$)
${}''E_2^{p, q} = \check{H}^p(K, \underline{H}^q(\mathcal{F}))$.
Puisque cette suite spectrale converge vers la cohomologie de
$\text{Tot}(A^{\bullet, \bullet})$, la démonstration est terminée.
\end{proof}

\begin{lemma}
\label{lemma-cech-spectral-sequence-verdier}
Soit $\mathcal{C}$ un site admettant des égalisateurs et des produits fibrés.
Soit $K$ un hyperrecouvrement.
Soit $\mathcal{F}$ un faisceau abélien. Il existe une
suite spectrale $(E_r, d_r)_{r \geq 0}$ avec
$$
E_2^{p, q} = \check{H}^p(K, \underline{H}^q(\mathcal{F}))
$$
qui converge vers les groupes de cohomologie globale $H^{p + q}(\mathcal{F})$.
\end{lemma}

\begin{proof}
C'est un cas particulier du lemme \ref{lemma-cech-spectral-sequence-variant}.
\end{proof}







\section{Recouvrements d'hyperrecouvrements}
\label{section-covering}

\noindent
Voici quelques façons de construire des hyperrecouvrements.
Notons que, puisque la catégorie
$\text{SR}(\mathcal{C}, X)$ admet des produits fibrés,
la catégorie des objets simpliciaux
de $\text{SR}(\mathcal{C}, X)$ admet elle aussi des produits fibrés,
voir Simplicial, lemme \ref{simplicial-lemma-fibre-product}.

\begin{lemma}
\label{lemma-funny-gamma}
Soit $\mathcal{C}$ un site admettant des produits fibrés.
Soit $X$ un objet de $\mathcal{C}$.
Soient $K, L, M$ des objets simpliciaux de $\text{SR}(\mathcal{C}, X)$.
Soient $a : K \to L$, $b : M \to L$ des morphismes.
Supposons que
\begin{enumerate}
\item $K$ soit un hyperrecouvrement de $X$ ;
\item le morphisme $M_0 \to L_0$ soit un recouvrement ;
\item pour tout $n \geq 0$, dans le diagramme
$$
\xymatrix{
M_{n + 1} \ar[dd] \ar[rr] \ar[rd]^\gamma &
&
(\text{cosk}_n \text{sk}_n M)_{n + 1} \ar[dd] \\
&
L_{n + 1}
\times_{(\text{cosk}_n \text{sk}_n L)_{n + 1}}
(\text{cosk}_n \text{sk}_n M)_{n + 1}
\ar[ld] \ar[ru]
& \\
L_{n + 1} \ar[rr] & & (\text{cosk}_n \text{sk}_n L)_{n + 1}
}
$$
la flèche $\gamma$ soit un recouvrement.
\end{enumerate}
Alors le produit fibré $K \times_L M$ est un hyperrecouvrement de $X$.
\end{lemma}

\begin{proof}
Le morphisme $(K \times_L M)_0 = K_0 \times_{L_0} M_0 \to K_0$
est un changement de base d'un recouvrement par (2), donc un recouvrement, voir le
lemme \ref{lemma-covering-permanence}. Et $K_0 \to \{X \to X\}$
est un recouvrement par (1). Ainsi, $(K \times_L M)_0 \to \{X \to X\}$
est un recouvrement par le lemme \ref{lemma-covering-permanence}. Par conséquent,
$K \times_L M$ satisfait la première condition de la définition
\ref{definition-hypercovering}.

\medskip\noindent
Il reste à vérifier que
$$
K_{n + 1} \times_{L_{n + 1}} M_{n + 1} = (K \times_L M)_{n + 1}
\longrightarrow
(\text{cosk}_n \text{sk}_n (K \times_L M))_{n + 1}
$$
est un recouvrement pour tout $n \geq 0$. Abrégeons comme suit :
$A = (\text{cosk}_n \text{sk}_n K)_{n + 1}$,
$B = (\text{cosk}_n \text{sk}_n L)_{n + 1}$, et
$C = (\text{cosk}_n \text{sk}_n M)_{n + 1}$.
Le foncteur $\text{cosk}_n \text{sk}_n$ commute aux produits fibrés,
voir Simplicial, lemme \ref{simplicial-lemma-cosk-fibre-product}.
Ainsi, le membre de droite ci-dessus est égal à $A \times_B C$.
Considérons le diagramme commutatif suivant
$$
\xymatrix{
K_{n + 1} \times_{L_{n + 1}} M_{n + 1} \ar[r] \ar[d] &
M_{n + 1} \ar[d] \ar[rd]_\gamma \ar[rrd] &
& \\
K_{n + 1} \ar[r] \ar[rd] &
L_{n + 1} \ar[rrd] &
L_{n + 1} \times_B C \ar[l] \ar[r] &
C \ar[d] \\
&
A \ar[rr] &
&
B
}
$$
Ce diagramme montre que
$$
K_{n + 1} \times_{L_{n + 1}} M_{n + 1}
=
(K_{n + 1} \times_B C)
\times_{(L_{n + 1} \times_B C), \gamma}
M_{n + 1}
$$
Or, $K_{n + 1} \times_B C \to A \times_B C$
est un changement de base du recouvrement $K_{n + 1} \to A$
par le morphisme $A \times_B C \to A$ ; c'est donc un
recouvrement. Par l'hypothèse (3), le morphisme $\gamma$ est un recouvrement.
Par conséquent, le morphisme
$$
(K_{n + 1} \times_B C)
\times_{(L_{n + 1} \times_B C), \gamma}
M_{n + 1}
\longrightarrow
K_{n + 1} \times_B C
$$
est un recouvrement, en tant que changement de base d'un recouvrement.
Le lemme en résulte, puisqu'un composé de recouvrements
est un recouvrement.
\end{proof}

\begin{lemma}
\label{lemma-product-hypercoverings}
Soit $\mathcal{C}$ un site admettant des produits fibrés.
Soit $X$ un objet de $\mathcal{C}$.
Si $K, L$ sont des hyperrecouvrements de $X$, alors
$K \times L$ est un hyperrecouvrement de $X$.
\end{lemma}

\begin{proof}
On peut soit le vérifier directement, soit utiliser le
lemme \ref{lemma-funny-gamma} ci-dessus et vérifier que $L \to \{X \to X\}$
possède la propriété (3).
\end{proof}


\noindent
Soit $\mathcal{C}$ un site admettant des produits fibrés.
Soit $X$ un objet de $\mathcal{C}$.
Puisque la catégorie $\text{SR}(\mathcal{C}, X)$ admet des coproduits et des
limites finies, il est licite de parler des objets
$U \times K$ et $\Hom(U, K)$ pour certains ensembles simpliciaux $U$
(par exemple ceux qui n'ont qu'un nombre fini de simplexes non dégénérés)
et tout objet simplicial $K$ de $\text{SR}(\mathcal{C}, X)$.
Voir Simplicial, sections
\ref{simplicial-section-product-with-simplicial-sets} et
\ref{simplicial-section-hom-from-simplicial-sets}.

\begin{lemma}
\label{lemma-covering}
Soit $\mathcal{C}$ un site admettant des produits fibrés.
Soit $X$ un objet de $\mathcal{C}$.
Soit $K$ un hyperrecouvrement de $X$.
Soit $k \geq 0$ un entier.
Soit $u : Z \to K_k$ un recouvrement
dans $\text{SR}(\mathcal{C}, X)$.
Alors il existe un morphisme d'hyperrecouvrements
$f: L \to K$ tel que $L_k \to K_k$
se factorise par $u$.
\end{lemma}

\begin{proof}
Notons $Y = K_k$. Soit $C[k]$ l'ensemble cosimplicial défini dans
Simplicial, exemple \ref{simplicial-example-simplex-cosimplicial-set}.
Nous utiliserons la description de $\Hom(C[k], Y)$ et $\Hom(C[k], Z)$
donnée dans
Simplicial, lemme \ref{simplicial-lemma-morphism-into-product}.
Il existe un morphisme canonique
$K \to \Hom(C[k], Y)$ correspondant à $\text{id} : K_k = Y \to Y$.
Considérons le morphisme $\Hom(C[k], Z) \to \Hom(C[k], Y)$
qui, sur les termes de degré $n$, est le morphisme
$$
\prod\nolimits_{\alpha : [k] \to [n]} Z
\longrightarrow
\prod\nolimits_{\alpha : [k] \to [n]} Y
$$
obtenu en utilisant le morphisme donné $Z \to Y$ sur chaque facteur. Posons
$$
L = K \times_{\Hom(C[k], Y)} \Hom(C[k], Z).
$$
Le morphisme $L_k \to K_k$ s'insère dans un diagramme commutatif
$$
\xymatrix{
L_k \ar[r] \ar[d] &
\prod_{\alpha : [k] \to [k]} Z \ar[r]^-{\text{pr}_{\text{id}_{[k]}}} \ar[d] &
Z \ar[d] \\
K_k \ar[r] &
\prod_{\alpha : [k] \to [k]} Y \ar[r]^-{\text{pr}_{\text{id}_{[k]}}} &
Y
}
$$
Puisque le composé des deux flèches du bas est l'identité,
nous concluons que nous avons la factorisation voulue.

\medskip\noindent
Il reste à montrer que $L$ est un hyperrecouvrement de $X$.
Pour ce faire, nous utiliserons le lemme \ref{lemma-funny-gamma}.
La condition (1) est satisfaite par hypothèse.
Pour (2), le morphisme
$$
\Hom(C[k], Z)_0 \to \Hom(C[k], Y)_0
$$
est un recouvrement, car il est isomorphe à $Z \to Y$ puisqu'il
n'existe qu'un seul morphisme $[k] \to [0]$.

\medskip\noindent
Considérons la condition (3) pour $n = 0$. Puisque
$(\text{cosk}_0 T)_1 = T \times T$
(Simplicial, exemple \ref{simplicial-example-cosk0})
et puisque $\Hom(C[k], Z)_1 = \prod_{\alpha : [k] \to [1]} Z$,
nous obtenons le diagramme
$$
\xymatrix{
\prod\nolimits_{\alpha : [k] \to [1]} Z \ar[r] \ar[d] &
Z \times Z \ar[d] \\
\prod\nolimits_{\alpha : [k] \to [1]} Y \ar[r] &
Y \times Y
}
$$
où les flèches horizontales correspondent à la projection sur les facteurs
associés aux deux $\alpha$ non surjectifs. Ainsi, la flèche $\gamma$
est le morphisme
$$
\prod\nolimits_{\alpha : [k] \to [1]} Z
\longrightarrow
\prod\nolimits_{\alpha : [k] \to [1]\text{ non surjectif}} Z
\times
\prod\nolimits_{\alpha : [k] \to [1]\text{ surjectif}} Y
$$
qui est un produit de recouvrements, donc un recouvrement d'après le
lemme \ref{lemma-covering-permanence}.

\medskip\noindent
Considérons la condition (3) pour $n > 0$. Nous affirmons qu'il existe une
application injective $\tau : S' \to S$ d'ensembles finis telle que, pour tout
objet $T$ de $\text{SR}(\mathcal{C}, X)$, le morphisme
\begin{equation}
\label{equation-map}
\Hom(C[k], T)_{n + 1} \to
(\text{cosk}_n \text{sk}_n \Hom(C[k], T))_{n + 1}
\end{equation}
soit isomorphe à la projection $\prod_{s \in S} T \to \prod_{s' \in S'} T$
fonctoriellement en $T$. Si cela est vrai, nous voyons, en raisonnant comme au paragraphe
précédent, que la flèche $\gamma$ est le morphisme
$$
\prod\nolimits_{s \in S} Z
\longrightarrow
\prod\nolimits_{s \in S'} Z
\times
\prod\nolimits_{s \not\in \tau(S')} Y
$$
qui est un produit de recouvrements, donc un recouvrement d'après le
lemme \ref{lemma-covering-permanence}. Par construction, nous avons
$\Hom(C[k], T)_{n + 1} = \prod_{\alpha : [k] \to [n + 1]} T$
(voir Simplicial, lemme \ref{simplicial-lemma-morphism-into-product}).
Nous prenons donc $S = \text{Map}([k], [n + 1])$.
D'autre part, Simplicial, lemme \ref{simplicial-lemma-formula-limit},
fournit une description des points de
$(\text{cosk}_n \text{sk}_n \Hom(C[k], T))_{n + 1}$
comme suites $(f_0, \ldots, f_{n + 1})$ de points de $\Hom(C[k], T)_n$
satisfaisant $d^n_{j - 1} f_i = d^n_i f_j$ pour $0 \leq i < j \leq n + 1$.
Nous pouvons écrire $f_i = (f_{i, \alpha})$, où $f_{i, \alpha}$ est un point de $T$
et $\alpha \in \text{Map}([k], [n])$. Les conditions se traduisent par
$$
f_{i, \delta^n_{j - 1} \circ \beta} = f_{j, \delta_i^n \circ \beta}
$$
pour tous $0 \leq i < j \leq n + 1$ et $\beta : [k] \to [n - 1]$. Ainsi, nous
voyons que
$$
S' = \{0, \ldots, n + 1\} \times \text{Map}([k], [n]) / \sim
$$
où la relation d'équivalence est engendrée par les équivalences
$$
(i, \delta^n_{j - 1} \circ \beta) \sim (j, \delta_i^n \circ \beta)
$$
pour $0 \leq i < j \leq n + 1$ et $\beta : [k] \to [n - 1]$.
Un calcul (omis) montre que le morphisme (\ref{equation-map})
correspond à l'application $S' \to S$ qui envoie $(i, \alpha)$ sur
$\delta^{n + 1}_i \circ \alpha \in S$. (Le lecteur trouvera peut-être rassurant
de constater que cette application est bien définie par la partie (1) de
Simplicial, lemme \ref{simplicial-lemma-relations-face-degeneracy}.)
Pour terminer la démonstration, il suffit de montrer que, si
$\alpha, \alpha' : [k] \to [n]$ et $0 \leq i < j \leq n + 1$
sont tels que
$$
\delta^{n + 1}_i \circ \alpha = \delta^{n + 1}_j \circ \alpha'
$$
alors nous avons $\alpha = \delta^n_{j - 1} \circ \beta$
et $\alpha' = \delta_i^n \circ \beta$ pour un certain $\beta : [k] \to [n - 1]$.
C'est immédiat, et nous l'omettons.
\end{proof}

\begin{lemma}
\label{lemma-covering-sheaf}
Soit $\mathcal{C}$ un site admettant des produits fibrés.
Soit $X$ un objet de $\mathcal{C}$.
Soit $K$ un hyperrecouvrement de $X$.
Soit $n \geq 0$ un entier.
Soit $u : \mathcal{F} \to F(K_n)$ un morphisme
de préfaisceaux qui devient surjectif
après faisceautisation.
Alors il existe un morphisme d'hyperrecouvrements
$f: L \to K$ tel que $F(f_n) : F(L_n) \to F(K_n)$
se factorise par $u$.
\end{lemma}

\begin{proof}
Écrivons $K_n = \{U_i \to X\}_{i \in I}$.
Ainsi, l'application $u$ est un morphisme de préfaisceaux d'ensembles
$u : \mathcal{F} \to \amalg h_{u_i}$.
L'hypothèse sur $u$ signifie que, pour tout
$i \in I$, il existe un recouvrement $\{U_{ij} \to U_i\}_{j \in I_i}$
du site $\mathcal{C}$ et un morphisme de préfaisceaux
$t_{ij} : h_{U_{ij}} \to \mathcal{F}$ tels que
$u \circ t_{ij}$ soit l'application $h_{U_{ij}} \to h_{U_i}$
provenant du morphisme $U_{ij} \to U_i$.
Posons $J = \amalg_{i \in I} I_i$, et soit
$\alpha : J \to I$ l'application évidente.
Pour $j \in J$, notons $V_j = U_{\alpha(j)j}$. Posons
$Z = \{V_j \to X\}_{j \in J}$.
Enfin, considérons le morphisme
$u' : Z \to K_n$ donné par $\alpha : J \to I$
et par les morphismes $V_j = U_{\alpha(j)j} \to U_{\alpha(j)}$
ci-dessus. Il s'agit clairement d'un recouvrement dans la
catégorie $\text{SR}(\mathcal{C}, X)$ et, par
construction, $F(u') : F(Z) \to F(K_n)$ se factorise par $u$.
Le résultat découle donc du lemme \ref{lemma-covering} ci-dessus.
\end{proof}


\section{Ajout de simplexes}
\label{section-adding-simplices}

\noindent
Dans cette section, nous démontrons quelques lemmes techniques dont nous aurons besoin plus loin.
Soit $\mathcal{C}$ un site admettant des produits fibrés.
Soit $X$ un objet de $\mathcal{C}$.
Comme nous l'avons indiqué dans la section \ref{section-covering} ci-dessus,
les objets $U \times K$ et $\Hom(U, K)$
sont définis pour certains ensembles simpliciaux $U$
et tout objet simplicial $K$ de $\text{SR}(\mathcal{C}, X)$.
Voir Simplicial, sections
\ref{simplicial-section-product-with-simplicial-sets} et
\ref{simplicial-section-hom-from-simplicial-sets}.

\begin{lemma}
\label{lemma-one-more-simplex}
Soit $\mathcal{C}$ un site admettant des produits fibrés.
Soit $X$ un objet de $\mathcal{C}$.
Soit $K$ un hyperrecouvrement de $X$.
Soient $U \subset V$ des ensembles simpliciaux tels que $U_n, V_n$
soient finis et non vides pour tout $n$.
Supposons que $U$ n'ait qu'un nombre fini de simplexes non dégénérés.
Supposons que $n \geq 0$ et $x \in V_n$,
$x \not \in U_n$, soient tels que
\begin{enumerate}
\item $V_i = U_i$ pour $i < n$ ;
\item $V_n = U_n \cup \{x\}$,
\item tout $z \in V_j$, $z \not \in U_j$, pour $j > n$,
soit dégénéré.
\end{enumerate}
Alors le morphisme
$$
\Hom(V, K)_0
\longrightarrow
\Hom(U, K)_0
$$
de $\text{SR}(\mathcal{C}, X)$ est un recouvrement.
\end{lemma}

\begin{proof}
Si $n = 0$, il s'ensuit aisément que $V = U \amalg \Delta[0]$
(voir ci-dessous). Dans ce cas, $\Hom(V, K)_0 =
\Hom(U, K)_0 \times K_0$. Le résultat découle alors
du lemme \ref{lemma-covering-permanence}.

\medskip\noindent
Soit $a : \Delta[n] \to V$ le morphisme associé à $x$
comme dans Simplicial, lemme \ref{simplicial-lemma-simplex-map}.
Écrivons $\partial \Delta[n] = i_{(n-1)!} \text{sk}_{n - 1} \Delta[n]$
pour le $(n - 1)$-squelette de $\Delta[n]$.
Soit $b : \partial \Delta[n] \to U$ la restriction
de $a$ au $(n - 1)$-squelette de $\Delta[n]$. D'après
Simplicial, lemme \ref{simplicial-lemma-glue-simplex},
nous avons $V = U \amalg_{\partial \Delta[n]} \Delta[n]$. D'après
Simplicial, lemme
\ref{simplicial-lemma-hom-from-coprod}
nous obtenons que
$$
\xymatrix{
\Hom(V, K)_0 \ar[r] \ar[d] &
\Hom(U, K)_0 \ar[d] \\
\Hom(\Delta[n], K)_0 \ar[r] &
\Hom(\partial \Delta[n], K)_0
}
$$
est un carré cartésien. Il suffit donc de montrer que
la flèche horizontale du bas est un recouvrement. D'après
Simplicial, lemme \ref{simplicial-lemma-cosk-shriek},
cette flèche s'identifie à
$$
K_n \to (\text{cosk}_{n - 1} \text{sk}_{n - 1} K)_n
$$
et est donc un recouvrement par définition d'un hyperrecouvrement.
\end{proof}

\begin{lemma}
\label{lemma-add-simplices}
Soit $\mathcal{C}$ un site admettant des produits fibrés.
Soit $X$ un objet de $\mathcal{C}$.
Soit $K$ un hyperrecouvrement de $X$.
Soient $U \subset V$ des ensembles simpliciaux tels que $U_n, V_n$
soient finis et non vides pour tout $n$.
Supposons que $U$ et $V$ n'aient qu'un nombre fini de simplexes non dégénérés.
Alors le morphisme
$$
\Hom(V, K)_0
\longrightarrow
\Hom(U, K)_0
$$
de $\text{SR}(\mathcal{C}, X)$ est un recouvrement.
\end{lemma}

\begin{proof}
D'après le lemme \ref{lemma-one-more-simplex}
ci-dessus, il suffit de démontrer un lemme élémentaire
sur les inclusions d'ensembles simpliciaux $U \subset V$ comme dans le
lemme. C'est exactement le résultat de
Simplicial, lemme \ref{simplicial-lemma-add-simplices}.
\end{proof}

\begin{lemma}
\label{lemma-degeneracy-maps-coverings}
Soit $\mathcal{C}$ un site admettant des produits fibrés. Soit $X$ un objet de
$\mathcal{C}$. Soit $K$ un hyperrecouvrement de $X$. Alors
\begin{enumerate}
\item $K_n$ est un recouvrement de $X$ pour tout $n \geq 0$ ;
\item $d^n_i : K_n \to K_{n - 1}$ est un recouvrement pour tous $n \geq 1$
et $0 \leq i \leq n$.
\end{enumerate}
\end{lemma}

\begin{proof}
Rappelons que $K_0$ est un recouvrement de $X$ par la
définition \ref{definition-hypercovering}
et que cela revient à dire que
$K_0 \to \{X \to X\}$ est un recouvrement au sens
de la définition \ref{definition-covering-SR}.
Ainsi, (1) découle de (2), car celui-ci montrera que le
composé
$K_n \to K_{n - 1} \to \ldots \to K_0 \to \{X \to X\}$
est un recouvrement par le lemme \ref{lemma-covering-permanence}.

\medskip\noindent
Démonstration de (2). Observons que
$\Mor(\Delta[n], K)_0 = K_n$ d'après
Simplicial, lemme \ref{simplicial-lemma-exists-hom-from-simplicial-set-finite}.
Par conséquent, (2) découle du lemme \ref{lemma-add-simplices}
appliqué aux $n + 1$ inclusions distinctes $\Delta[n - 1] \to \Delta[n]$.
\end{proof}

\begin{remark}
\label{remark-P-covering}
Un cas particulier utile des lemmes \ref{lemma-add-simplices} et
\ref{lemma-degeneracy-maps-coverings} est le suivant.
Supposons que nous ayons une catégorie $\mathcal{C}$ admettant des produits fibrés.
Soit $P \subset \text{Arrows}(\mathcal{C})$ une partie
stable par changement de base, stable par composition
et contenant tous les isomorphismes. On appelle alors
{\it $P$-hyperrecouvrement} une augmentation $a : U \to X$
issue d'un objet simplicial de $\mathcal{C}$ telle que
\begin{enumerate}
\item $U_0 \to X$ appartienne à $P$ ;
\item $U_1 \to U_0 \times_X U_0$ appartienne à $P$ ;
\item $U_{n + 1} \to (\text{cosk}_n\text{sk}_n U)_{n + 1}$
appartienne à $P$ pour $n \geq 1$.
\end{enumerate}
La catégorie $\mathcal{C}/X$ admet toutes les limites finies ; les
cosquelettes utilisés dans la formulation ci-dessus existent donc
(voir Catégories, lemme \ref{categories-lemma-finite-limits-exist}).
Nous affirmons alors que les morphismes $U_n \to X$ et $d^n_i : U_n \to U_{n - 1}$
appartiennent à $P$. Cela découle des lemmes
précédents en faisant de $\mathcal{C}$ un site dont les recouvrements
sont les $\{f : V \to U\}$ tels que $f \in P$, et en prenant $K$ donné par
$K_n = \{U_n \to X\}$.
\end{remark}


\section{Homotopies}
\label{section-homotopies}

\noindent
Soit $\mathcal{C}$ un site admettant des produits fibrés.
Soit $X$ un objet de $\mathcal{C}$.
Soit $L$ un objet simplicial de $\text{SR}(\mathcal{C}, X)$.
D'après
Simplicial, lemme \ref{simplicial-lemma-exists-hom-from-simplicial-set-finite},
il existe un objet $\Hom(\Delta[1], L)$
de la catégorie $\text{Simp}(\text{SR}(\mathcal{C}, X))$ qui représente le
foncteur
$$
T
\longmapsto
\Mor_{\text{Simp}(\text{SR}(\mathcal{C}, X))}(\Delta[1] \times T, L)
$$
Il existe un morphisme canonique
$$
\Hom(\Delta[1], L) \to L \times L
$$
provenant de $e_i : \Delta[0] \to \Delta[1]$ et de l'identification
$\Hom(\Delta[0], L) = L$.

\begin{lemma}
\label{lemma-hom-hypercovering}
Soit $\mathcal{C}$ un site admettant des produits fibrés.
Soit $X$ un objet de $\mathcal{C}$.
Soit $L$ un objet simplicial de $\text{SR}(\mathcal{C}, X)$.
Soit $n \geq 0$. Considérons le diagramme commutatif
\begin{equation}
\label{equation-diagram}
\xymatrix{
\Hom(\Delta[1], L)_{n + 1} \ar[r] \ar[d] &
(\text{cosk}_n \text{sk}_n \Hom(\Delta[1], L))_{n + 1} \ar[d] \\
(L \times L)_{n + 1} \ar[r] &
(\text{cosk}_n \text{sk}_n (L \times L))_{n + 1}
}
\end{equation}
provenant du morphisme défini ci-dessus.
Nous pouvons identifier les termes de ce diagramme comme suit,
où
$\partial \Delta[n + 1] = i_{n!}\text{sk}_n \Delta[n + 1]$
est le $n$-squelette du $(n + 1)$-simplexe :
\begin{eqnarray*}
\Hom(\Delta[1], L)_{n + 1}
& = &
\Hom(\Delta[1] \times \Delta[n + 1], L)_0 \\
(\text{cosk}_n \text{sk}_n \Hom(\Delta[1], L))_{n + 1}
& = &
\Hom(\Delta[1] \times \partial \Delta[n + 1], L)_0 \\
(L \times L)_{n + 1}
& = &
\Hom(
(\Delta[n + 1] \amalg \Delta[n + 1], L)_0 \\
(\text{cosk}_n \text{sk}_n (L \times L))_{n + 1}
& = &
\Hom(
\partial \Delta[n + 1]
\amalg
\partial \Delta[n + 1], L)_0
\end{eqnarray*}
et les morphismes entre ces objets de $\text{SR}(\mathcal{C}, X)$
proviennent du diagramme commutatif d'ensembles simpliciaux
\begin{equation}
\label{equation-dual-diagram}
\xymatrix{
\Delta[1] \times \Delta[n + 1] &
\Delta[1] \times \partial\Delta[n + 1] \ar[l] \\
\Delta[n + 1] \amalg \Delta[n + 1] \ar[u] &
\partial\Delta[n + 1] \amalg \partial\Delta[n + 1]
\ar[l] \ar[u]
}
\end{equation}
En outre, le produit fibré de la flèche du bas et de la
flèche de droite dans (\ref{equation-diagram}) est égal à
$$
\Hom(U, L)_0
$$
où $U \subset \Delta[1] \times \Delta[n + 1]$
est le plus petit sous-ensemble simplicial contenant les images de
$\Delta[n + 1] \amalg \Delta[n + 1]$ et de
$\Delta[1] \times \partial\Delta[n + 1]$.
\end{lemma}

\begin{proof}
La première et la troisième égalité sont données par
Simplicial, lemme \ref{simplicial-lemma-exists-hom-from-simplicial-set-finite}.
La deuxième et la quatrième résultent de ce lemme combiné à
Simplicial, lemme \ref{simplicial-lemma-cosk-shriek}.
La dernière assertion découle du fait que
$U$ est la somme amalgamée des flèches du bas et de droite du
diagramme (\ref{equation-dual-diagram}), par
Simplicial, lemme \ref{simplicial-lemma-hom-from-coprod}.
Pour voir que $U$ est égal à cette somme amalgamée, il suffit
de constater que l'intersection de
$\Delta[n + 1] \amalg \Delta[n + 1]$ et de
$\Delta[1] \times \partial\Delta[n + 1]$
dans $\Delta[1] \times \Delta[n + 1]$ est égale à
$\partial\Delta[n + 1] \amalg \partial\Delta[n + 1]$.
Nous laissons cela au lecteur.
\end{proof}

\begin{lemma}
\label{lemma-homotopy}
Soit $\mathcal{C}$ un site admettant des produits fibrés.
Soit $X$ un objet de $\mathcal{C}$.
Soient $K, L$ des hyperrecouvrements de $X$.
Soient $a, b : K \to L$ des morphismes d'hyperrecouvrements.
Il existe un morphisme d'hyperrecouvrements
$c : K' \to K$ tel que $a \circ c$ soit homotope
à $b \circ c$.
\end{lemma}

\begin{proof}
Considérons le diagramme commutatif suivant
$$
\xymatrix{
K' \ar@{=}[r]^-{def} \ar[rd]_c &
K \times_{(L \times L)} \Hom(\Delta[1], L)
\ar[r] \ar[d] & \Hom(\Delta[1], L) \ar[d] \\
& K \ar[r]^{(a, b)} & L \times L
}
$$
Par la propriété fonctorielle de $\Hom(\Delta[1], L)$,
le composé des morphismes horizontaux
correspond à un morphisme $K' \times \Delta[1] \to L$ qui
définit une homotopie entre $c \circ a$ et $c \circ b$.
Ainsi, si nous pouvons montrer que $K'$ est un
hyperrecouvrement de $X$, nous obtenons le lemme.
Pour ce faire, nous appliquerons le lemme \ref{lemma-funny-gamma}
à la paire de morphismes $K \to L \times L$
et $\Hom(\Delta[1], L) \to L \times L$.
La condition (1) du lemme \ref{lemma-funny-gamma} est satisfaite.
La condition (2) du lemme \ref{lemma-funny-gamma} est vérifiée car
$\Hom(\Delta[1], L)_0 = L_1$, et le morphisme
$(d^1_0, d^1_1) : L_1 \to L_0 \times L_0$ est un
recouvrement de $\text{SR}(\mathcal{C}, X)$ par notre
hypothèse selon laquelle $L$ est un hyperrecouvrement.
Pour démontrer la condition (3) du lemme \ref{lemma-funny-gamma},
nous utilisons le lemme \ref{lemma-hom-hypercovering} ci-dessus. D'après
ce lemme, le morphisme $\gamma$ de la condition (3) du lemme
\ref{lemma-funny-gamma} est le morphisme
$$
\Hom(\Delta[1] \times \Delta[n + 1], L)_0
\longrightarrow
\Hom(U, L)_0
$$
où $U \subset \Delta[1] \times \Delta[n + 1]$.
D'après le lemme \ref{lemma-add-simplices},
c'est un recouvrement ; l'assertion est donc démontrée.
\end{proof}

\begin{remark}
\label{remark-contractible-category}
Notons que le point essentiel de la démonstration est l'emploi du
lemme \ref{lemma-add-simplices}. Ce lemme
est tout à fait général et ne dépend pas de la
forme exacte des ensembles simpliciaux (pourvu qu'ils
n'aient qu'un nombre fini de simplexes non dégénérés).
Il paraît tout à fait raisonnable de s'attendre à un résultat
du type suivant :
étant donné un morphisme quelconque $a : K \times \partial \Delta[k]
\to L$, où $K$ et $L$ sont des hyperrecouvrements, il
existe un morphisme d'hyperrecouvrements $c : K' \to K$
et un morphisme $g : K' \times \Delta[k] \to L$
tels que
$g|_{K' \times \partial \Delta[k]} =
a \circ (c \times \text{id}_{\partial \Delta[k]})$.
Autrement dit, la catégorie des hyperrecouvrements est, en un
sens convenable, contractile.
\end{remark}

















\section{Cohomologie et hyperrecouvrements}
\label{section-cohomology}

\noindent
Soit $\mathcal{C}$ un site admettant des produits fibrés.
Soit $X$ un objet de $\mathcal{C}$.
Soit $\mathcal{F}$ un faisceau de groupes abéliens sur $\mathcal{C}$.
Soient $K, L$ des hyperrecouvrements de $X$.
Si $a, b : K \to L$ sont des applications homotopes,
alors $\mathcal{F}(a), \mathcal{F}(b) : \mathcal{F}(K) \to \mathcal{F}(L)$
sont des applications homotopes, voir
Simplicial, lemme \ref{simplicial-lemma-functorial-homotopy}.
Elles ont donc le même effet sur les groupes de cohomologie des
complexes de cochaînes associés, voir
Simplicial, lemme \ref{simplicial-lemma-homotopy-s-Q}.
Nous allons utiliser ce fait pour définir la limite inductive sur tous les
hyperrecouvrements.

\medskip\noindent
Notons provisoirement $\text{HC}(\mathcal{C}, X)$
la catégorie dont les objets sont les hyperrecouvrements de $X$
et dont les morphismes sont les applications entre hyperrecouvrements de $X$
à homotopie près. Nous avons vu qu'il
s'agit d'une catégorie et non d'une « grande » catégorie,
voir le lemme \ref{lemma-hypercoverings-set}.
La catégorie opposée à $\text{HC}(\mathcal{C}, X)$
sera la catégorie d'indices de notre diagramme, voir
Catégories, section \ref{categories-section-limits} pour la terminologie.
Considérons le diagramme
$$
\check{H}^i(-, \mathcal{F}) :
\text{HC}(\mathcal{C}, X)^{opp}
\longrightarrow
\textit{Ab}.
$$
D'après les lemmes \ref{lemma-product-hypercoverings} et \ref{lemma-homotopy}
et la remarque sur les homotopies ci-dessus, ce diagramme est filtrant, voir
Catégories, définition \ref{categories-definition-directed}.
Ainsi, la limite inductive
$$
\check{H}^i_{\text{HC}}(X, \mathcal{F}) =
\colim_{K \in \text{HC}(\mathcal{C}, X)} \check{H}^i(K, \mathcal{F})
$$
admet une description particulièrement simple (voir la référence citée).

\begin{theorem}
\label{theorem-cohomology-hypercoverings}
Soit $\mathcal{C}$ un site admettant des produits fibrés.
Soit $X$ un objet de $\mathcal{C}$. Soit $i \geq 0$.
Les foncteurs
\begin{eqnarray*}
\textit{Ab}(\mathcal{C}) & \longrightarrow & \textit{Ab} \\
\mathcal{F} & \longmapsto & H^i(X, \mathcal{F}) \\
\mathcal{F} & \longmapsto & \check{H}^i_{\text{HC}}(X, \mathcal{F})
\end{eqnarray*}
sont canoniquement isomorphes.
\end{theorem}

\begin{proof}[Démonstration utilisant les suites spectrales.]
Supposons que $\xi \in H^p(X, \mathcal{F})$ pour un certain $p \geq 0$.
Montrons que $\xi$ appartient à l'image de l'application
$\check{H}^p(X, \mathcal{F}) \to H^p(X, \mathcal{F})$ du
lemme \ref{lemma-cech-spectral-sequence}
pour un certain hyperrecouvrement $K$ de $X$.

\medskip\noindent
Cela est vrai si $p = 0$ par le lemme \ref{lemma-h0-cech}.
Si $p = 1$, choisissons un hyperrecouvrement de {\v C}ech $K$ de $X$ comme dans
l'exemple \ref{example-cech}, en partant d'un recouvrement
$K_0 = \{U_i \to X\}$ du site $\mathcal{C}$ tel que
$\xi|_{U_i} = 0$, voir
Cohomologie sur les sites,
lemme \ref{sites-cohomology-lemma-kill-cohomology-class-on-covering}.
Il résulte immédiatement de la suite spectrale
du lemme \ref{lemma-cech-spectral-sequence} que, dans ce cas, $\xi$ provient
d'un élément de $\check{H}^1(K, \mathcal{F})$.
En général, choisissons un hyperrecouvrement quelconque $K$ de $X$ tel
que $\xi$ s'envoie sur zéro dans $\underline{H}^p(\mathcal{F})(K_0)$
(en utilisant l'exemple \ref{example-cech} et
Cohomologie sur les sites,
lemme \ref{sites-cohomology-lemma-kill-cohomology-class-on-covering},
une fois encore).
D'après la suite spectrale du lemme \ref{lemma-cech-spectral-sequence},
l'obstruction à ce que $\xi$ provienne d'un élément de
$\check{H}^p(K, \mathcal{F})$ est une suite d'éléments
$\xi_1, \ldots, \xi_{p - 1}$ tels que
$\xi_q \in \check{H}^{p - q}(K, \underline{H}^q(\mathcal{F}))$
(plus précisément, les images des $\xi_q$ dans certains sous-quotients
de ces groupes).

\medskip\noindent
Nous pouvons remplacer par récurrence l'hyperrecouvrement $K$ par des raffinements
sur lesquels les obstructions $\xi_1, \ldots, \xi_{p - 1}$ se restreignent à zéro
(et non seulement leurs images
dans les sous-quotients : il n'y a donc ici aucune subtilité). Supposons en effet être
déjà parvenus à la situation où
$\xi_{q + 1}, \ldots, \xi_{p - 1}$ sont nuls.
Notons que $\xi_q \in \check{H}^{p - q}(K, \underline{H}^q(\mathcal{F}))$
est la classe d'un certain élément
$$
\tilde \xi_q \in
\underline{H}^q(\mathcal{F})(K_{p - q}) =
\prod H^q(U_i, \mathcal{F})
$$
si $K_{p - q} = \{U_i \to X\}_{i \in I}$. Soit $\xi_{q, i}$
la composante de $\tilde \xi_q$ dans $H^q(U_i, \mathcal{F})$.
Comme $q \geq 1$, nous pouvons utiliser
Cohomologie sur les sites,
lemme \ref{sites-cohomology-lemma-kill-cohomology-class-on-covering},
une fois encore, pour choisir des recouvrements $\{U_{i, j} \to U_i\}$
du site tels que chaque restriction $\xi_{q, i}|_{U_{i, j}} = 0$.
Considérons l'objet $Z = \{U_{i, j} \to X\}$ de la catégorie
$\text{SR}(\mathcal{C}, X)$ et son morphisme évident
$u : Z \to K_{p - q}$. Il est clair que $u$ est un recouvrement, voir la
définition \ref{definition-covering-SR}. D'après le
lemme \ref{lemma-covering}, il
existe un morphisme $L \to K$ d'hyperrecouvrements de $X$ tel que
$L_{p - q} \to K_{p - q}$ se factorise par $u$. Alors, clairement,
l'image de $\xi_q$ dans $\underline{H}^q(\mathcal{F})(L_{p - q})$
est nulle. Puisque la suite spectrale du
lemme \ref{lemma-cech-spectral-sequence}
est fonctorielle, cela signifie qu'après avoir remplacé $K$ par $L$, nous atteignons la
situation où $\xi_q, \ldots, \xi_{p - 1}$ sont tous nuls.
En poursuivant ainsi, nous aboutissons à un hyperrecouvrement où ils sont tous
nuls ; par conséquent, $\xi$ appartient à l'image de l'application
$\check{H}^p(X, \mathcal{F}) \to H^p(X, \mathcal{F})$.

\medskip\noindent
Supposons que $K$ soit un hyperrecouvrement de $X$, que
$\xi \in \check{H}^p(K, \mathcal{F})$ et que l'image de
$\xi$ par l'application
$\check{H}^p(X, \mathcal{F}) \to H^p(X, \mathcal{F})$ du
lemme \ref{lemma-cech-spectral-sequence}
soit nulle. Pour terminer la démonstration du théorème, nous devons montrer
qu'il existe un morphisme d'hyperrecouvrements $L \to K$ tel que
$\xi$ se restreigne à zéro dans $\check{H}^p(L, \mathcal{F})$.
D'après la suite spectrale du lemme \ref{lemma-cech-spectral-sequence},
l'annulation de l'image de $\xi$ dans $H^p(X, \mathcal{F})$
signifie qu'il existe des éléments $\xi_1, \ldots, \xi_{p - 2}$
tels que $\xi_q \in \check{H}^{p - 1 - q}(K, \underline{H}^q(\mathcal{F}))$
(plus précisément, leurs images dans certains sous-quotients)
et tels que la somme des images $d_{q + 1}^{p - 1 - q, q}\xi_q$ (dans la suite
spectrale) soit égale à $\xi$. Par exactement le même mécanisme que ci-dessus,
nous pouvons donc trouver un morphisme d'hyperrecouvrements $L \to K$ tel que
les restrictions des éléments $\xi_q$, $q = 1, \ldots, p - 2$,
dans $\check{H}^{p - 1 - q}(L, \underline{H}^q(\mathcal{F}))$ soient nulles.
Il s'ensuit que $\xi$ est nul puisque le morphisme $L \to K$
induit un morphisme de suites spectrales d'après le
lemme \ref{lemma-cech-spectral-sequence}.
\end{proof}

\begin{proof}[Démonstration sans utiliser les suites spectrales.]
Nous avons vu le résultat pour $i = 0$, voir le lemme \ref{lemma-h0-cech}.
Nous savons que les foncteurs $H^i(X, -)$ forment un $\delta$-foncteur universel, voir
Catégories dérivées, lemme \ref{derived-lemma-higher-derived-functors}.
Pour démontrer le théorème, il suffit de montrer que
la suite de foncteurs $\check{H}^i_{HC}(X, -)$ forme un
$\delta$-foncteur. En effet, nous savons que la cohomologie de {\v C}ech
est nulle en degrés strictement positifs sur les faisceaux injectifs (lemme \ref{lemma-injective-trivial-cech}),
et nous pouvons alors appliquer
Homologie, lemme \ref{homology-lemma-efface-implies-universal}.

\medskip\noindent
Soit
$$
0 \to \mathcal{F} \to \mathcal{G} \to \mathcal{H} \to 0
$$
une suite exacte courte de faisceaux abéliens sur $\mathcal{C}$.
Soit $\xi \in \check{H}^p_{HC}(X, \mathcal{H})$. Choisissons un hyperrecouvrement
$K$ de $X$ et un élément $\sigma \in \mathcal{H}(K_p)$ représentant
$\xi$ en cohomologie. Il lui correspond une suite exacte de
complexes
$$
0 \to s(\mathcal{F}(K)) \to s(\mathcal{G}(K)) \to s(\mathcal{H}(K))
$$
mais rien ne garantit qu'il y ait aussi un zéro à droite, et c'est
la seule chose qui
nous empêche de définir $\delta(\xi)$ par une simple application du
lemme du serpent. Rappelons que
$$
\mathcal{H}(K_p) = \prod \mathcal{H}(U_i)
$$
si $K_p = \{U_i \to X\}$. Soit $\sigma =\prod \sigma_i$ avec
$\sigma_i \in \mathcal{H}(U_i)$. Puisque $\mathcal{G} \to \mathcal{H}$ est
une surjection de faisceaux, nous voyons qu'il existe des recouvrements
$\{U_{i, j} \to U_i\}$ tels que $\sigma_i|_{U_{i, j}}$ soit
l'image d'un certain élément $\tau_{i, j} \in \mathcal{G}(U_{i, j})$.
Considérons l'objet $Z = \{U_{i, j} \to X\}$ de la catégorie
$\text{SR}(\mathcal{C}, X)$ et son morphisme évident
$u : Z \to K_p$. Il est clair que $u$ est un recouvrement, voir la
définition \ref{definition-covering-SR}. D'après le
lemme \ref{lemma-covering}, il
existe un morphisme $L \to K$ d'hyperrecouvrements de $X$ tel que
$L_p \to K_p$ se factorise par $u$. Après avoir remplacé $K$ par $L$,
nous pouvons donc supposer que $\sigma$ est l'image d'un
élément $\tau \in \mathcal{G}(K_p)$. Notons que $d(\sigma) = 0$,
mais pas nécessairement $d(\tau) = 0$. Ainsi, $d(\tau) \in \mathcal{F}(K_{p + 1})$
est un cocycle. Dans cette situation, nous définissons
$\delta(\xi)$ comme la classe du cocycle $d(\tau)$ dans
$\check{H}^{p + 1}_{HC}(X, \mathcal{F})$.

\medskip\noindent
À ce stade, plusieurs points sont à vérifier :
(a) $\delta(\xi)$ ne dépend pas du choix de $\tau$ ;
(b) $\delta(\xi)$ ne dépend pas du choix de l'hyperrecouvrement
$L \to K$ sur lequel $\sigma$ se relève ;
(c) $\delta(\xi)$ ne dépend ni de l'hyperrecouvrement initial ni de
$\sigma$ choisi pour représenter $\xi$. Nous omettons la vérification de
(a), (b) et (c) ; l'indépendance du choix des hyperrecouvrements
se ramène en fait aux lemmes \ref{lemma-product-hypercoverings}
et \ref{lemma-homotopy}. Nous omettons aussi de vérifier que
$\delta$ est fonctoriel relativement aux morphismes de suites exactes
courtes de faisceaux abéliens sur $\mathcal{C}$.

\medskip\noindent
Enfin, nous devons vérifier qu'avec cette définition de $\delta$,
notre suite exacte courte de faisceaux abéliens ci-dessus donne une
suite exacte longue de groupes de cohomologie de {\v C}ech.
Montrons d'abord que, si $\delta(\xi) = 0$ (avec $\xi$ comme ci-dessus), alors
$\xi$ est l'image d'un certain élément
$\xi' \in \check{H}^p_{HC}(X, \mathcal{G})$.
En effet, si $\delta(\xi) = 0$, alors, avec les notations ci-dessus, nous
voyons que la classe de $d(\tau)$ est nulle dans
$\check{H}^{p + 1}_{HC}(X, \mathcal{F})$. Il existe donc
un morphisme d'hyperrecouvrements $L \to K$ tel que la restriction
de $d(\tau)$ à un élément de $\mathcal{F}(L_{p + 1})$ soit
égale à $d(\upsilon)$ pour un certain $\upsilon \in \mathcal{F}(L_p)$.
Cela implique que $\tau|_{L_p} + \upsilon$ forme un
cocycle et détermine une classe $\xi' \in \check{H}^p(L, \mathcal{G})$
qui s'envoie sur $\xi$, comme voulu.

\medskip\noindent
Nous omettons la démonstration du fait que, si $\xi' \in \check{H}^{p + 1}_{HC}(X, \mathcal{F})$
s'envoie sur zéro dans $\check{H}^{p + 1}_{HC}(X, \mathcal{G})$, alors il est
égal à $\delta(\xi)$ pour un certain $\xi \in \check{H}^p_{HC}(X, \mathcal{H})$.
\end{proof}

\noindent
Nous déduisons maintenant le cas de Verdier du
théorème \ref{theorem-cohomology-hypercoverings}
par un artifice.

\begin{proposition}
\label{proposition-cohomology-hypercoverings}
Soit $\mathcal{C}$ un site admettant des produits fibrés et les produits de deux objets.
Soit $\mathcal{F}$ un faisceau abélien sur $\mathcal{C}$.
Soit $i \geq 0$. Alors
\begin{enumerate}
\item pour tout $\xi \in H^i(\mathcal{F})$, il existe un hyperrecouvrement
$K$ tel que $\xi$ appartienne à l'image de l'application canonique
$\check{H}^i(K, \mathcal{F}) \to H^i(\mathcal{F})$ ;
\item si $K, L$ sont des hyperrecouvrements et si $\xi_K \in \check{H}^i(K, \mathcal{F})$,
$\xi_L \in \check{H}^i(L, \mathcal{F})$ sont des éléments qui s'envoient
sur le même élément de $H^i(\mathcal{F})$, alors il existe
un hyperrecouvrement $M$ et des morphismes $M \to K$ et $M \to L$ tels
que $\xi_K$ et $\xi_L$ s'envoient sur le même élément de
$\check{H}^i(M, \mathcal{F})$.
\end{enumerate}
Autrement dit, modulo les questions ensemblistes, les groupes de
cohomologie de $\mathcal{F}$ sur $\mathcal{C}$ sont la limite inductive des
groupes de cohomologie de {\v C}ech de $\mathcal{F}$ sur tous les hyperrecouvrements.
\end{proposition}

\begin{proof}
Ce résultat est une conséquence immédiate du
théorème \ref{theorem-cohomology-hypercoverings}.
En effet, nous pouvons remplacer artificiellement $\mathcal{C}$ par un site
légèrement plus grand $\mathcal{C}'$ tel que
(I) $\mathcal{C}'$ admette un objet final $X$ et (II)
les hyperrecouvrements dans $\mathcal{C}$ soient plus ou moins la
même chose que les hyperrecouvrements de $X$ dans $\mathcal{C}'$.
Mais cette construction exige un travail assez important de
vérification des détails.

\medskip\noindent
Appelons une famille de morphismes $\{U_i \to U\}$ de $\mathcal{C}$
de but fixé un {\it recouvrement faible} si le morphisme de faisceaux associé à
l'application $\coprod_{i \in I} h_{U_i} \to h_U$ est surjectif.
Construisons un nouveau site $\mathcal{C}'$ comme suit :
\begin{enumerate}
\item comme catégorie, posons $\Ob(\mathcal{C}') = \Ob(\mathcal{C}) \amalg \{X\}$
et ajoutons un unique morphisme vers $X$ depuis tout objet de $\mathcal{C}'$ ;
\item $\mathcal{C}'$ admet des produits fibrés puisque les produits fibrés et les produits
de deux objets existent dans $\mathcal{C}$ ;
\item les recouvrements de $\mathcal{C}'$ sont les recouvrements faibles de $\mathcal{C}$,
ainsi que les $\{U_i \to X\}_{i \in I}$ tels que soit $U_i = X$
pour un certain $i$, soit $U_i \not = X$ pour tout $i$ et l'application
$\coprod h_{U_i} \to *$ de préfaisceaux sur $\mathcal{C}$ devienne
surjective après faisceautisation sur $\mathcal{C}$ ;
\item nous appliquons Ensembles, lemme \ref{sets-lemma-coverings-site},
pour restreindre les recouvrements et obtenir notre site $\mathcal{C}'$.
\end{enumerate}
Alors $\Sh(\mathcal{C}') = \Sh(\mathcal{C})$, car le foncteur d'inclusion
$\mathcal{C} \to \mathcal{C}'$ est un foncteur cocontinu spécial
(voir Sites, définition \ref{sites-definition-special-cocontinuous-functor}).
Nous omettons les vérifications immédiates.

\medskip\noindent
Choisissons un recouvrement $\{U_i \to X\}$ de $\mathcal{C}'$ tel que $U_i$ soit un
objet de $\mathcal{C}$ pour tout $i$ (c'est possible car
$\mathcal{C} \to \mathcal{C}'$ est cocontinu spécial).
Alors $K_0 = \{U_i \to X\}$ est un recouvrement du
site $\mathcal{C}'$ construit ci-dessus. Considérons $K_0$ comme un objet de
$\text{SR}(\mathcal{C}', X)$ et posons $K_{init} = \text{cosk}_0(K_0)$.
Alors $K_{init}$ est un hyperrecouvrement de $X$, voir
l'exemple \ref{example-cech}. Notons que tout $K_{init, n}$ est de la forme
$\{W_j \to X\}$ avec $W_j \in \Ob(\mathcal{C})$.

\medskip\noindent
Démonstration de (1). Choisissons $\xi \in H^i(\mathcal{F}) = H^i(X, \mathcal{F}')$,
où $\mathcal{F}'$ est le faisceau abélien sur $\mathcal{C}'$ correspondant
à $\mathcal{F}$ sur $\mathcal{C}$. D'après le
théorème \ref{theorem-cohomology-hypercoverings},
il existe un morphisme d'hyperrecouvrements $K' \to K_{init}$
de $X$ dans $\mathcal{C}'$ tel que $\xi$ provienne d'un élément
de $\check{H}^i(K', \mathcal{F})$.
Écrivons $K'_n = \{U_{n, j} \to X\}$. Puisque $K'_n$ s'envoie dans
$K_{init, n}$, nous voyons que $U_{n, j}$ est un objet de $\mathcal{C}$.
Nous pouvons donc définir un objet simplicial $K$ de $\text{SR}(\mathcal{C})$
en posant $K_n = \{U_{n, j}\}$. Puisque les recouvrements de
$\mathcal{C}'$ constitués de familles de morphismes de $\mathcal{C}$
sont des recouvrements faibles, nous voyons que $K$ est un hyperrecouvrement au sens
de la définition \ref{definition-hypercovering-variant}.
Enfin, puisque $\mathcal{F}'$ est l'unique faisceau sur $\mathcal{C}'$
dont la restriction à $\mathcal{C}$ est égale à $\mathcal{F}$,
nous voyons que les complexes de {\v C}ech $s(\mathcal{F}(K))$
et $s(\mathcal{F}'(K'))$ sont identiques, d'où (1).
(Compatibilité avec l'application vers les groupes de cohomologie omise.)

\medskip\noindent
Démonstration de (2). Soient $K$ et $L$ des hyperrecouvrements dans $\mathcal{C}$.
Soient $K'$ et $L'$ les objets simpliciaux de $\text{SR}(\mathcal{C}', X)$
obtenus à partir de $K$ et $L$ par le foncteur
$\text{SR}(\mathcal{C}) \to \text{SR}(\mathcal{C}', X)$,
$\{U_i\} \mapsto \{U_i \to X\}$. Comme précédemment, les complexes de
{\v C}ech sont égaux ; nous obtenons donc $\xi_{K'}$ et
$\xi_{L'}$ qui s'envoient sur la même classe de cohomologie de $\mathcal{F}'$
sur $\mathcal{C}'$. Après avoir éventuellement élargi notre choix
de recouvrements dans $\mathcal{C}'$ (en raison d'une question ensembliste),
nous pouvons supposer que $K'$ et $L'$ sont des hyperrecouvrements de $X$ dans
$\mathcal{C}'$ ; cela résulte de notre définition des hyperrecouvrements dans la
définition \ref{definition-hypercovering-variant} et
du fait que les recouvrements faibles dans $\mathcal{C}$ donnent des recouvrements dans
$\mathcal{C}'$. D'après le
théorème \ref{theorem-cohomology-hypercoverings},
il existe un hyperrecouvrement $M'$ de $X$ dans $\mathcal{C}'$
et des morphismes $M' \to K'$, $M' \to L'$ et $M' \to K_{init}$
tels que $\xi_{K'}$ et $\xi_{L'}$ se restreignent au même élément de
$\check{H}^i(M', \mathcal{F})$. En explicitant cet énoncé comme ci-dessus,
nous trouvons que (2) est vrai.
\end{proof}





\section{Hyperrecouvrements d'espaces}
\label{section-hypercoverings-spaces}

\noindent
La théorie précédente présente un certain intérêt même dans le cas des espaces
topologiques. Dans ce cas, nous pouvons expliciter ce qu'est un hyperrecouvrement et voir
ce que dit réellement le résultat.

\medskip\noindent
Soit $X$ un espace topologique. Considérons le site $X_{Zar}$ de
Sites, exemple \ref{sites-example-site-topological}. Rappelons
qu'un objet de $X_{Zar}$ est simplement un ouvert de $X$ et que les morphismes
de $X_{Zar}$ correspondent simplement aux inclusions. Qu'est donc un
hyperrecouvrement de $X$ pour le site $X_{Zar}$ ?

\medskip\noindent
Explicitons d'abord la définition \ref{definition-SR}.
Un objet de $\text{SR}(X_{Zar}, X)$ est simplement donné par un ensemble
$I$ et, pour chaque $i \in I$, un ouvert $U_i \subset X$.
Notons-le $\{U_i\}_{i \in I}$, puisqu'il ne peut y avoir de
confusion sur le morphisme $U_i \to X$.
Un morphisme $\{U_i\}_{i \in I} \to \{V_j\}_{j \in J}$
entre deux tels objets est donné par une application d'ensembles
$\alpha : I \to J$ telle que $U_i \subset V_{\alpha(i)}$ pour tout
$i \in I$. Quand un tel morphisme est-il un recouvrement ? C'est le cas
si et seulement si, pour tout $j \in J$, nous avons
$V_j = \bigcup_{i\in I, \ \alpha(i) = j} U_i$ (et c'est
un recouvrement du site $X_{Zar}$).

\medskip\noindent
À l'aide de ce qui précède, nous obtenons la description suivante d'un hyperrecouvrement
dans le site $X_{Zar}$. Un hyperrecouvrement de $X$ dans $X_{Zar}$
est donné par les données suivantes :
\begin{enumerate}
\item un ensemble simplicial $I$ (voir
Simplicial, section \ref{simplicial-section-simplicial-set}) ;
\item pour chaque $n \geq 0$ et tout $i \in I_n$, un ouvert $U_i \subset X$.
\end{enumerate}
Nous noterons une telle collection de données $(I, \{U_i\})$.
Pour qu'il s'agisse d'un hyperrecouvrement de $X$, nous exigeons
les propriétés suivantes :
\begin{itemize}
\item pour $i \in I_n$ et $0 \leq a \leq n$,
nous avons $U_i \subset U_{d^n_a(i)}$ ;
\item pour $i \in I_n$ et $0 \leq a \leq n$, nous avons $U_i = U_{s^n_a(i)}$ ;
\item nous avons
\begin{equation}
\label{equation-covering-X}
X = \bigcup\nolimits_{i \in I_0} U_i,
\end{equation}
\item pour tous $i_0, i_1 \in I_0$, nous avons
\begin{equation}
\label{equation-covering-two}
U_{i_0} \cap U_{i_1} =
\bigcup\nolimits_{i \in I_1, \ d^1_0(i) = i_0, \ d^1_1(i) = i_1} U_i,
\end{equation}
\item pour tout $n \geq 1$ et tout
$(i_0, \ldots, i_{n + 1}) \in (I_n)^{n + 2}$ tel que
$d^n_{b - 1}(i_a) = d^n_a(i_b)$ pour tous $0\leq a < b\leq n + 1$,
nous avons
\begin{equation}
\label{equation-covering-general}
U_{i_0} \cap \ldots \cap U_{i_{n + 1}} =
\bigcup\nolimits_{i \in I_{n + 1},
\ d^{n + 1}_a(i) = i_a, \ a = 0, \ldots, n + 1} U_i,
\end{equation}
\item chacun des recouvrements ouverts (\ref{equation-covering-X}),
(\ref{equation-covering-two}) et (\ref{equation-covering-general})
est un élément de $\text{Cov}(X_{Zar})$
(il s'agit d'une condition ensembliste qui borne
la taille des ensembles d'indices des recouvrements).
\end{itemize}
Les conditions (\ref{equation-covering-X}) et
(\ref{equation-covering-two}) devraient être familières, par exemple à la lecture du
chapitre sur les faisceaux sur les espaces, et la condition
(\ref{equation-covering-general}) en est la généralisation naturelle.

\begin{remark}
\label{remark-not-covering-set}
Une particularité de cette description est que, si l'une des intersections multiples
$U_{i_0} \cap \ldots \cap U_{i_{n + 1}}$ est vide, alors
le recouvrement du membre de droite peut être le recouvrement vide.
Ainsi, les applications
$I_{n + 1} \to (\text{cosk}_n\text{sk}_n I)_{n + 1}$ ne sont pas nécessairement surjectives.
Cela signifie que la réalisation géométrique de $I$ peut être un espace
intéressant (non contractile).

\medskip\noindent
En fait, soit $I'_n \subset I_n$ la partie
constituée des simplexes $i \in I_n$ tels que
$U_i \not = \emptyset$. Il est facile de voir que $I' \subset I$
est un sous-ensemble simplicial et que $(I', \{U_i\})$ est un hyperrecouvrement.
Nous pouvons donc toujours raffiner un hyperrecouvrement en un hyperrecouvrement dans lequel
aucun des ouverts $U_i$ n'est vide.
\end{remark}

\begin{remark}
\label{remark-repackage-into-simplicial-space}
Reformulons ces informations d'une autre manière encore.
Supposons que $(I, \{U_i\})$ soit un hyperrecouvrement de
l'espace topologique $X$. À partir de ces données, nous pouvons construire
un espace topologique simplicial $U_\bullet$ en posant
$$
U_n = \coprod\nolimits_{i \in I_n} U_i,
$$
et, pour $\varphi : [n] \to [m]$ donné, en prenant pour
morphisme $U(\varphi) : U_n \to U_m$ celui
provenant des inclusions $U_i \subset U_{\varphi(i)}$
pour $i \in I_n$. Cet espace topologique simplicial est muni
d'une augmentation $\epsilon : U_\bullet \to X$.
Avec ce morphisme, l'espace simplicial $U_\bullet$ devient
un hyperrecouvrement de $X$ le long duquel on a la descente cohomologique
au sens de \cite[Expos\'e Vbis]{SGA4}.
Autrement dit, $H^n(U_\bullet, \epsilon^*\mathcal{F}) = H^n(X, \mathcal{F})$.
(Insérer ici une référence ultérieure à la cohomologie sur les espaces
simpliciaux et à la descente cohomologique formulée en ces termes.)
Supposons que $\mathcal{F}$ soit un faisceau abélien sur $X$.
Dans ce cas, la suite spectrale du lemme \ref{lemma-cech-spectral-sequence}
devient la suite spectrale de terme $E_1$
$$
E_1^{p, q} = H^q(U_p, \epsilon_q^*\mathcal{F})
\Rightarrow
H^{p + q}(U_\bullet, \epsilon^*\mathcal{F}) = H^{p + q}(X, \mathcal{F})
$$
qui compare la cohomologie totale de $\epsilon^*\mathcal{F}$
aux groupes de cohomologie de $\mathcal{F}$ sur les morceaux
de $U_\bullet$. (Insérer ici une référence ultérieure à cette suite
spectrale.)
\end{remark}

\noindent
En topologie, nous voulons souvent trouver des hyperrecouvrements de $X$ tels
que tous les $U_i$ proviennent d'une base donnée de la topologie
de $X$ et que tous les recouvrements
(\ref{equation-covering-two}) et (\ref{equation-covering-general})
appartiennent à une collection cofinale donnée de recouvrements.
Voici deux lemmes illustrant cela.

\begin{lemma}
\label{lemma-basis-hypercovering}
Soit $X$ un espace topologique.
Soit $\mathcal{B}$ une base de la topologie de $X$.
Il existe un hyperrecouvrement $(I, \{U_i\})$ de $X$
tel que chaque $U_i$ appartienne à $\mathcal{B}$.
\end{lemma}

\begin{proof}
Soit $n \geq 0$.
Appelons {\it hyperrecouvrement $n$-tronqué} de $X$ une donnée
constituée d'un ensemble simplicial $n$-tronqué $I$ et, pour chaque
$i \in I_a$, $0 \leq a \leq n$, d'un ouvert $U_i$ de $X$, telle que
les conditions définissant un hyperrecouvrement soient satisfaites chaque fois qu'elles ont un sens.
Autrement dit, nous n'imposons les relations d'inclusion et les conditions de
recouvrement que lorsque tous les simplexes qui y interviennent
sont des $a$-simplexes avec $a \leq n$. Le lemme en découlera si nous pouvons montrer
que, étant donné un hyperrecouvrement $n$-tronqué $(I, \{U_i\})$ avec
tous les $U_i \in \mathcal{B}$, nous pouvons l'étendre en un hyperrecouvrement
$(n + 1)$-tronqué sans ajouter de $a$-simplexe pour $a \leq n$.
Procédons comme suit. Considérons d'abord l'ensemble simplicial
$(n + 1)$-tronqué $I'$ défini par
$I' = \text{sk}_{n + 1}(\text{cosk}_n I)$.
Rappelons que
$$
I'_{n + 1} =
\left\{
\begin{matrix}
(i_0, \ldots, i_{n + 1}) \in (I_n)^{n + 2} \text{ tels que}\\
d^n_{b - 1}(i_a) = d^n_a(i_b) \text{ pour tous }0\leq a < b\leq n + 1
\end{matrix}
\right\}
$$
Si $i' \in I'_{n + 1}$ est dégénéré, disons $i' = s^n_a(i)$, nous posons
$U_{i'} = U_i$ (cela nous est de toute façon imposé par la deuxième condition).
Nous posons aussi $J_{i'} = \{i'\}$ dans ce cas.
Si $i' \in I'_{n + 1}$ est non dégénéré, disons
$i' = (i_0, \ldots, i_{n + 1})$, nous choisissons un ensemble
$J_{i'}$ et un recouvrement ouvert
\begin{equation}
\label{equation-choose-covering}
U_{i_0} \cap \ldots \cap U_{i_{n + 1}} =
\bigcup\nolimits_{i \in J_{i'}} U_i,
\end{equation}
avec $U_i \in \mathcal{B}$ pour $i \in J_{i'}$.
Posons
$$
I_{n + 1} = \coprod\nolimits_{i' \in I'_{n + 1}} J_{i'}
$$
Il existe une application canonique $\pi : I_{n + 1} \to I'_{n + 1}$ qui est
une bijection au-dessus de l'ensemble des simplexes dégénérés de $I'_{n + 1}$, par
construction.
Pour $i \in I_{n + 1}$, nous définissons $d^{n + 1}_a(i) = d^{n + 1}_a(\pi(i))$.
Pour $i \in I_n$, nous définissons $s^n_a(i) \in I_{n + 1}$ comme l'unique
simplexe situé au-dessus du simplexe dégénéré $s^n_a(i) \in I'_{n + 1}$.
Nous omettons de vérifier que cela définit un hyperrecouvrement
$(n + 1)$-tronqué de $X$.
\end{proof}

\begin{lemma}
\label{lemma-quasi-separated-quasi-compact-hypercovering}
Soit $X$ un espace topologique.
Soit $\mathcal{B}$ une base de la topologie de $X$.
Supposons que
\begin{enumerate}
\item $X$ soit quasi-compact ;
\item tout $U \in \mathcal{B}$ soit un ouvert quasi-compact ;
\item l'intersection de deux ouverts quasi-compacts quelconques de
$X$ soit quasi-compacte.
\end{enumerate}
Alors il existe un hyperrecouvrement $(I, \{U_i\})$ de $X$ ayant les
propriétés suivantes :
\begin{enumerate}
\item chaque $U_i$ appartient à la base $\mathcal{B}$ ;
\item chacun des $I_n$ est un ensemble fini ; en particulier,
\item chacun des recouvrements (\ref{equation-covering-X}),
(\ref{equation-covering-two}) et (\ref{equation-covering-general})
est fini.
\end{enumerate}
\end{lemma}

\begin{proof}
Cela découle directement de la construction donnée dans la démonstration du
lemme \ref{lemma-basis-hypercovering} si nous choisissons des recouvrements finis
par des éléments de $\mathcal{B}$ dans (\ref{equation-choose-covering}).
Détails omis.
\end{proof}







\section{Construction d'hyperrecouvrements}
\label{section-hypercovering-sites}

\noindent
Soit $\mathcal{C}$ un site. Dans cette section, nous considérerons un
objet simplicial de $\text{SR}(\mathcal{C})$ comme suit.
Comme d'habitude, nous posons $K_n = K([n])$ et notons $K(\varphi) : K_n \to K_m$
le morphisme associé à $\varphi : [m] \to [n]$.
Nous pouvons écrire $K_n = \{U_{n, i}\}_{i \in I_n}$. Pour
$\varphi : [m] \to [n]$, le morphisme $K(\varphi) : K_n \to K_m$
est donné par une application $\alpha(\varphi) : I_n \to I_m$ et des morphismes
$f_{\varphi, i} : U_{n, i} \to U_{m, \alpha(\varphi)(i)}$
pour $i \in I_n$. Le fait que $K$ soit un objet simplicial de
$\text{SR}(\mathcal{C})$ implique que $(I_n, \alpha(\varphi))$
est un ensemble simplicial
et que $f_{\psi, \alpha(\varphi)(i)} \circ f_{\varphi, i} =
f_{\varphi \circ \psi, i}$ lorsque $\psi : [l] \to [m]$.

\begin{lemma}
\label{lemma-split}
Soit $\mathcal{C}$ un site. Soit $K$ un objet simplicial $r$-tronqué
de $\text{SR}(\mathcal{C})$. Les conditions suivantes sont équivalentes :
\begin{enumerate}
\item $K$ est scindé (Simplicial, définition \ref{simplicial-definition-split}) ;
\item $f_{\varphi, i} : U_{n, i} \to U_{m, \alpha(\varphi)(i)}$
est un isomorphisme pour $r \geq n \geq 0$,
$\varphi : [m] \to [n]$ surjective et $i \in I_n$ ;
\item $f_{\sigma^n_j, i} : U_{n, i} \to U_{n + 1, \alpha(\sigma^n_j)(i)}$
est un isomorphisme pour $0 \leq j \leq n < r$, $i \in I_n$.
\end{enumerate}
Il en va de même pour les objets simpliciaux si, dans (2) et (3),
nous posons $r = \infty$.
\end{lemma}

\begin{proof}
La scission d'un ensemble simplicial est unique et est donnée par
les indices non dégénérés $N(I_n)$ en chaque degré $n$, voir
Simplicial, lemme \ref{simplicial-lemma-splitting-simplicial-sets}.
Le coproduit de deux objets $\{U_i\}_{i \in I}$ et $\{U_j\}_{j \in J}$
de $\text{SR}(\mathcal{C})$ est donné par $\{U_l\}_{l \in I \amalg J}$
avec les notations évidentes. Une scission de $K$ doit donc être donnée par
$N(K_n) = \{U_i\}_{i \in N(I_n)}$. L'équivalence de (1) et (2)
résulte alors de l'explicitation des définitions. L'équivalence de (2)
et (3) découle du fait que toute surjection
$\varphi : [m] \to [n]$ est un composé de morphismes
$\sigma^k_j$ avec $k = n, n + 1, \ldots, m - 1$.
\end{proof}

\begin{lemma}
\label{lemma-hypercovering-object}
Soit $\mathcal{C}$ un site admettant des produits fibrés.
Soit $\mathcal{B} \subset \Ob(\mathcal{C})$ une partie.
Supposons que
\begin{enumerate}
\item tout objet $U$ de $\mathcal{C}$ admette un recouvrement
$\{U_j \to U\}_{j \in J}$ avec $U_j \in \mathcal{B}$ ;
\item si $\{U_j \to U\}_{j \in J}$ est un recouvrement
avec $U_j \in \mathcal{B}$ et si $\{U' \to U\}$ est un morphisme avec
$U' \in \mathcal{B}$, alors $\{U_j \to U\}_{j \in J} \amalg \{U' \to U\}$
soit un recouvrement.
\end{enumerate}
Alors, pour tout $X$ dans $\mathcal{C}$, il existe un hyperrecouvrement $K$
de $X$ tel que $K_n = \{U_{n, i}\}_{i \in I_n}$
avec $U_{n, i} \in \mathcal{B}$ pour tout $i \in I_n$.
\end{lemma}

\begin{proof}
La démonstration du lemme \ref{lemma-basis-hypercovering}
constitue une préparation à celle-ci, et
nous invitons le lecteur à la lire d'abord.

\medskip\noindent
Remplaçons d'abord $\mathcal{C}$ par le site $\mathcal{C}/X$.
Nous pouvons alors supposer que $X$ est l'objet final
de $\mathcal{C}$ et que $\mathcal{C}$ admet toutes les limites finies
(Catégories, lemme \ref{categories-lemma-finite-limits-exist}).

\medskip\noindent
Soit $n \geq 0$. Appelons
{\it hyperrecouvrement $n$-tronqué de type $\mathcal{B}$ de $X$}
une donnée constituée d'un objet simplicial $n$-tronqué $K$
de $\text{SR}(\mathcal{C})$
tel que, pour $i \in I_a$, $0 \leq a \leq n$,
nous ayons $U_{a, i} \in \mathcal{B}$, que
$K_0$ soit un recouvrement de $X$ et que
$K_{a + 1} \to (\text{cosk}_a \text{sk}_a K)_{a + 1}$
pour $a = 0, \ldots, n - 1$
soit un recouvrement comme dans la définition \ref{definition-covering-SR}.

\medskip\noindent
Puisque $X$ admet un recouvrement $\{U_{0, i} \to X\}_{i \in I_0}$
avec $U_i \in \mathcal{B}$ par hypothèse, nous obtenons un hyperrecouvrement
$0$-tronqué de type $\mathcal{B}$ de $X$. Observons que tout hyperrecouvrement
$0$-tronqué de type $\mathcal{B}$ de $X$ est scindé, voir le
lemme \ref{lemma-split}.

\medskip\noindent
Le lemme en découlera si nous pouvons montrer, pour $n \geq 0$, que, étant donné un
hyperrecouvrement scindé $n$-tronqué de type $\mathcal{B}$, noté $K$, de $X$,
nous pouvons l'étendre en un
hyperrecouvrement scindé $(n + 1)$-tronqué de type $\mathcal{B}$ de $X$.

\medskip\noindent
Construction de l'extension. Considérons l'objet simplicial $(n + 1)$-tronqué
$K' = \text{sk}_{n + 1}(\text{cosk}_n K)$ de $\text{SR}(\mathcal{C})$.
Écrivons
$$
K'_{n + 1} = \{U'_{n + 1, i}\}_{i \in I'_{n + 1}}
$$
Puisque $K = \text{sk}_n K'$, nous avons $K_a = K'_a$ pour $0 \leq a \leq n$.
Pour tout $i' \in I'_{n + 1}$, nous choisissons un recouvrement
\begin{equation}
\label{equation-choose-covering-B}
\{g_{n + 1, j} : U_{n + 1, j} \to U'_{n + 1, i'}\}_{j \in J_{i'}}
\end{equation}
avec $U_{n + 1, j} \in \mathcal{B}$ pour $j \in J_{i'}$.
C'est possible par notre hypothèse sur $\mathcal{B}$ dans le lemme.
Pour $0 \leq m \leq n$, notons $N_m \subset I_m$ la partie des
indices non dégénérés. Nous posons
$$
I_{n + 1} =
\coprod\nolimits_{\varphi : [n + 1] \to [m]\text{ surjective, }0\leq m \leq n}
N_m \amalg
\coprod\nolimits_{i' \in I'_{n + 1}} J_{i'}
$$
Pour $j \in I_{n + 1}$, nous posons
$$
U_{n + 1, j} =
\left\{
\begin{matrix}
U_{m, i} & \text{si} &
j = (\varphi, i) & \text{où} & \varphi : [n + 1] \to [m], i \in N_m \\
U_{n + 1, j} & \text{si} &
j \in J_{i'} & \text{où} & i' \in I'_{n + 1}
\end{matrix}
\right.
$$
avec les notations évidentes. Posons $K_{n + 1} = \{U_{n + 1, j}\}_{j \in I_{n + 1}}$.
Par construction, $U_{n + 1, j}$ appartient
à $\mathcal{B}$ pour tout $j \in I_{n + 1}$. Définissons des applications
compatibles
$$
I_{n + 1} \to I'_{n + 1}
\quad\text{et}\quad
K_{n + 1} \to K'_{n + 1}
$$
Plus précisément, la première application est donnée par
$(\varphi, i) \mapsto \alpha'(\varphi)(i)$ et
$(j \in J_{i'}) \mapsto i'$.
Pour la seconde application, nous utilisons les morphismes
$$
f'_{\varphi, i} : U_{m, i} \to U'_{n + 1, \alpha'(\varphi)(i)}
\quad\text{et}\quad
g_{n + 1, j} : U_{n + 1, j} \to U'_{n + 1, i'}
$$
Nous affirmons que le morphisme
$$
K_{n + 1} \to K'_{n + 1} =
(\text{cosk}_n \text{sk}_n K')_{n + 1} =
(\text{cosk}_n K)_{n + 1}
$$
est un recouvrement comme dans la définition \ref{definition-covering-SR}.
En effet, si $i' \in I'_{n + 1}$, alors soit $i'$ est non dégénéré
et l'image réciproque de $i'$ dans $I_{n + 1}$ est égale à $J_{i'}$,
auquel cas nous obtenons un recouvrement de $U'_{n + 1, i'}$ par notre choix
(\ref{equation-choose-covering-B}), soit $i'$ est dégénéré et
l'image réciproque de $i'$ dans $I_{n + 1}$ est
$J_{i'} \amalg \{(\varphi, i)\}$ pour une unique paire $(\varphi, i)$,
auquel cas nous obtenons un recouvrement par notre choix (\ref{equation-choose-covering-B})
et l'hypothèse (2) du lemme.

\medskip\noindent
Pour terminer la démonstration, nous devons définir les morphismes
$K(\varphi) : K_{n + 1} \to K_m$ correspondant aux morphismes
$\varphi : [m] \to [n + 1]$, $0 \leq m \leq n$, et les morphismes
$K(\varphi) : K_m \to K_{n + 1}$ correspondant aux morphismes
$\varphi : [n + 1] \to [m]$, $0 \leq m \leq n$,
de sorte que les relations de composition convenables soient satisfaites.
Pour ceux du premier type, nous utilisons le composé
$$
K_{n + 1} \to K'_{n + 1} \xrightarrow{K'(\varphi)} K'_m = K_m
$$
pour définir $K(\varphi) : K_{n + 1} \to K_m$.
Pour ceux du second type, supposons donné $\varphi : [n + 1] \to [m]$,
$0 \leq m \leq n$. Nous définissons le morphisme correspondant
$K(\varphi) : K_m \to K_{n + 1}$ comme suit :
\begin{enumerate}
\item pour $i \in I_m$, il existe une unique application surjective
$\psi : [m] \to [m_0]$ et un unique $i_0 \in I_{m_0}$ non dégénéré
tels que $\alpha(\psi)(i_0) = i$\footnote{Par exemple, si $i$ est
non dégénéré, alors $m = m_0$ et $\psi = \text{id}_{[m]}$.} ;
\item nous posons $\varphi_0 = \psi_0 \circ \varphi : [n + 1] \to [m_0]$
et envoyons
$i \in I_m$ sur $(\varphi_0, i_0) \in I_{n + 1}$ ; autrement dit,
$\alpha(\varphi)(i) = (\varphi_0, i_0)$ ;
\item le morphisme
$f_{\varphi, i} : U_{m, i} \to U_{n + 1, \alpha(\varphi)(i)} = U_{m_0, i_0}$
est l'inverse de l'isomorphisme $f_{\psi, i_0} : U_{m_0, i_0} \to U_{m, i}$
(voir le lemme \ref{lemma-split}).
\end{enumerate}
Nous omettons la vérification immédiate mais fastidieuse que cela définit
un hyperrecouvrement scindé $(n + 1)$-tronqué de type $\mathcal{B}$ de $X$
qui prolonge l'hyperrecouvrement $n$-tronqué donné. En fait, tout est clair
d'après ce qui précède, sauf la vérification que les morphismes
$K(\varphi)$ se composent correctement pour tout $\varphi : [a] \to [b]$
avec $0 \leq a, b \leq n + 1$.
\end{proof}

\begin{lemma}
\label{lemma-hypercovering-site}
Soit $\mathcal{C}$ un site admettant des égalisateurs et des produits fibrés.
Soit $\mathcal{B} \subset \Ob(\mathcal{C})$ une partie. Supposons
que tout objet de $\mathcal{C}$ admette un recouvrement
dont les membres appartiennent à $\mathcal{B}$.
Alors il existe un hyperrecouvrement $K$ tel que
$K_n = \{U_i\}_{i \in I_n}$ avec $U_i \in \mathcal{B}$
pour tout $i \in I_n$.
\end{lemma}

\begin{proof}
Cette démonstration est presque identique à celle du
lemme \ref{lemma-hypercovering-object}. Nous
n'expliquerons que les différences.

\medskip\noindent
Soit $n \geq 1$. Appelons
{\it hyperrecouvrement $n$-tronqué de type $\mathcal{B}$}
une donnée constituée d'un objet simplicial
$n$-tronqué $K$ de $\text{SR}(\mathcal{C})$
tel que, pour $i \in I_a$, $0 \leq a \leq n$,
nous ayons $U_{a, i} \in \mathcal{B}$ et que
\begin{enumerate}
\item $F(K_0)^\# \to *$ soit surjectif ;
\item $F(K_1)^\# \to F(K_0)^\# \times F(K_0)^\#$ soit surjectif ;
\item $F(K_{a + 1})^\# \to F((\text{cosk}_a \text{sk}_a K)_{a + 1})^\#$
soit surjectif pour $a = 1, \ldots, n - 1$.
\end{enumerate}
Construisons d'abord explicitement un hyperrecouvrement scindé $1$-tronqué de type $\mathcal{B}$.

\medskip\noindent
Prenons $I_0 = \mathcal{B}$ et $K_0 = \{U\}_{U \in \mathcal{B}}$.
Alors (1) vaut par notre hypothèse sur $\mathcal{B}$. Posons
$$
\Omega =
\{(U, V, W, a, b) \mid U, V, W \in \mathcal{B}, a : U \to V, b : U \to W\}
$$
Posons ensuite $I_1 = I_0 \amalg \Omega$. Pour $i \in I_1$, posons
$U_{1, i} = U_{0, i}$ si $i \in I_0$ et $U_{1, i} = U$
si $i = (U, V, W, a, b) \in \Omega$. L'application
$K(\sigma^0_0) : K_0 \to K_1$ correspond à
l'inclusion $\alpha(\sigma^0_0) : I_0 \to I_1$
et à l'identité $f_{\sigma^0_0, i} : U_{0, i} \to U_{1, i}$
sur les objets. Les applications $K(\delta^1_0), K(\delta^1_1) : K_1 \to K_0$
correspondent aux deux applications $I_1 \to I_0$ qui sont
l'identité sur $I_0 \subset I_1$ et envoient $(U, V, W, a, b) \in \Omega \subset I_1$
sur $V$, respectivement sur $W$. Les morphismes correspondants
$f_{\delta^1_0, i}, f_{\delta^1_1, i} : U_{1, i} \to U_{0, i}$ sont
l'identité si $i \in I_0$ et $a, b$ si $i = (U, V, W, a, b) \in \Omega$.
La raison pour laquelle (2) vaut est que toute section de
$F(K_0)^\# \times F(K_0)^\#$ au-dessus d'un objet $U$ de $\mathcal{C}$
provient, après avoir remplacé $U$ par les membres d'un recouvrement,
d'une application $U \to F(K_0) \times F(K_0)$.
Cela signifie à son tour que nous avons $V, W \in \mathcal{B}$
et deux morphismes $U \to V$ et $U \to W$. En remplaçant encore
$U$ par les membres d'un recouvrement, nous pouvons supposer $U \in \mathcal{B}$,
comme voulu.

\medskip\noindent
Le lemme en découlera si nous pouvons montrer que, étant donné un
hyperrecouvrement scindé $n$-tronqué de type $\mathcal{B}$, noté $K$, pour $n \geq 1$,
nous pouvons l'étendre en un hyperrecouvrement scindé $(n + 1)$-tronqué de type $\mathcal{B}$.
Ici, l'argument suit exactement la démonstration du
lemme \ref{lemma-hypercovering-object}.
Nous omettons les détails précis, à l'exception des remarques suivantes.
Premièrement, nous n'avons pas besoin de l'hypothèse (2) dans la démonstration du présent
lemme, car il n'est pas nécessaire que le morphisme
$K_{n + 1} \to (\text{cosk}_n K)_{n + 1}$ soit un recouvrement ;
il suffit qu'il induise une surjection sur les faisceaux d'ensembles associés,
ce qui découle de
Sites, lemme \ref{sites-lemma-covering-surjective-after-sheafification}.
Deuxièmement, l'hypothèse selon laquelle $\mathcal{C}$ admet des produits fibrés et des égalisateurs
garantit que $\text{SR}(\mathcal{C})$ admet des produits fibrés
et des égalisateurs, et que $F$ commute à ceux-ci
(lemme \ref{lemma-coprod-prod-SR}). Cela suffit à
garantir l'existence des foncteurs cosquelette utilisés (voir
Simplicial, remarque \ref{simplicial-remark-existence-cosk} et
Catégories, lemme \ref{categories-lemma-fibre-products-equalizers-exist}).
\end{proof}

\begin{lemma}
\label{lemma-hypercovering-morphism-sites}
Soit $f : \mathcal{C} \to \mathcal{D}$ un morphisme de sites
défini par le foncteur $u : \mathcal{D} \to \mathcal{C}$.
Supposons que $\mathcal{D}$ et $\mathcal{C}$ admettent des égalisateurs et
des produits fibrés, et que $u$ commute à ceux-ci.
Si un objet simplicial $K$ de $\text{SR}(\mathcal{D})$
est un hyperrecouvrement, alors $u(K)$ est un hyperrecouvrement.
\end{lemma}

\begin{proof}
Si nous écrivons $K_n = \{U_{n, i}\}_{i \in I_n}$ comme dans l'introduction
de cette section, alors $u(K)$ est l'objet de $\text{SR}(\mathcal{C})$
défini par $u(K_n) = \{u(U_i)\}_{i \in I_n}$.
D'après Sites, lemme \ref{sites-lemma-pullback-representable-sheaf},
nous avons $f^{-1}h_U^\# = h_{u(U)}^\#$ pour $U \in \Ob(\mathcal{D})$.
Cela signifie que $f^{-1}F(K_n)^\# = F(u(K_n))^\#$ pour tout $n$.
Vérifions les conditions (1), (2), (3) pour que $u(K)$ soit un
hyperrecouvrement suivant la définition \ref{definition-hypercovering-variant}.
Puisque $f^{-1}$ est un foncteur exact, nous obtenons le morphisme
$$
F(u(K_0))^\# = f^{-1}F(K_0)^\# \to f^{-1}* = *
$$
qui est surjectif comme image inverse d'un morphisme surjectif, d'où (1).
De même,
$$
F(u(K_1))^\# = f^{-1}F(K_1)^\# \to
f^{-1} (F(K_0) \times F(K_0))^\# = F(u(K_0))^\# \times F(u(K_0))^\#
$$
est surjectif comme image inverse, d'où (2). Pour la condition (3),
afin de conclure par la même méthode, il suffit d'avoir
$$
F((\text{cosk}_n \text{sk}_n u(K))_{n + 1})^\# =
f^{-1}F((\text{cosk}_n \text{sk}_n K)_{n + 1})^\#
$$
Ce qui précède montre que $f^{-1}F(-) = F(u(-))$. Il suffit donc de montrer
que $u$ commute aux limites utilisées pour définir
$(\text{cosk}_n \text{sk}_n K)_{n + 1}$ pour $n \geq 1$.
D'après Simplicial, remarque \ref{simplicial-remark-existence-cosk},
ces limites sont des limites finies connexes auxquelles $u$ commute
par hypothèse.
\end{proof}

\begin{lemma}
\label{lemma-hypercovering-continuous-functor}
Soient $\mathcal{C}$ et $\mathcal{D}$ des sites. Soit
$u : \mathcal{D} \to \mathcal{C}$ un foncteur continu.
Supposons que $\mathcal{D}$ et $\mathcal{C}$ admettent des produits fibrés
et que $u$ commute à ceux-ci. Soient $Y \in \mathcal{D}$ et
$K \in \text{SR}(\mathcal{D}, Y)$ un hyperrecouvrement de $Y$.
Alors $u(K)$ est un hyperrecouvrement de $u(Y)$.
\end{lemma}

\begin{proof}
Cela est plus facile que la démonstration du lemme \ref{lemma-hypercovering-morphism-sites},
car la notion d'hyperrecouvrement d'un objet est plus forte, voir les
définitions \ref{definition-hypercovering} et \ref{definition-covering-SR}.
En effet, $u$ envoie les recouvrements sur des recouvrements par définition
d'un morphisme de sites. Il suffit de vérifier que $u$ commute aux
limites utilisées pour définir
$(\text{cosk}_n \text{sk}_n K)_{n + 1}$ pour $n \geq 1$.
C'est clair, car le foncteur induit
$\mathcal{D}/Y \to \mathcal{C}/X$ commute à toutes les limites finies
(et les catégories source et but admettent toutes les limites finies d'après
Catégories, lemme \ref{categories-lemma-finite-limits-exist}).
\end{proof}

\begin{lemma}
\label{lemma-w-contractible}
Soit $\mathcal{C}$ un site. Soit $\mathcal{B} \subset \Ob(\mathcal{C})$
une partie. Supposons que
\begin{enumerate}
\item $\mathcal{C}$ admette des produits fibrés ;
\item pour tout $X \in \Ob(\mathcal{C})$, il existe un recouvrement fini
$\{U_i \to X\}_{i \in I}$ avec $U_i \in \mathcal{B}$ ;
\item si $\{U_i \to X\}_{i \in I}$ est un recouvrement fini avec
$U_i \in \mathcal{B}$ et si $U \to X$ est un morphisme avec $U \in \mathcal{B}$,
alors $\{U_i \to X\}_{i \in I} \amalg \{U \to X\}$ soit un recouvrement.
\end{enumerate}
Alors, pour tout $X$, il existe un hyperrecouvrement $K$ de $X$
tel que chaque $K_n = \{U_{n, i} \to X\}_{i \in I_n}$ vérifie
que $I_n$ est fini et $U_{n, i} \in \mathcal{B}$.
\end{lemma}

\begin{proof}
Ce lemme est l'analogue du
lemme \ref{lemma-quasi-separated-quasi-compact-hypercovering}
pour les sites. Pour le démontrer, nous suivons exactement la démonstration du
lemme \ref{lemma-hypercovering-object}
en prêtant attention aux deux points suivants :
\begin{enumerate}
\item[(a)] nous choisissons notre recouvrement initial $\{U_{0, i} \to X\}_{i \in I_0}$
avec $U_{0, i} \in \mathcal{B}$ de sorte que l'ensemble d'indices $I_0$ soit fini ;
\item[(b)] lors du choix des recouvrements
(\ref{equation-choose-covering-B})
nous choisissons $J_{i'}$ fini.
\end{enumerate}
Le lecteur voit aisément qu'avec ces modifications nous obtenons
des ensembles d'indices $I_n$ finis pour tout $n$.
\end{proof}

\begin{remark}
\label{remark-taking-disjoint-unions}
Soit $\mathcal{C}$ un site. Soient
$K$ et $L$ des objets de $\text{SR}(\mathcal{C})$.
Écrivons $K = \{U_i\}_{i \in I}$ et $L = \{V_j\}_{j \in J}$.
Supposons que $U = \coprod_{i \in I} U_i$ et $V = \coprod_{j \in J} V_j$
existent. Nous obtenons alors
$$
\Mor_{\text{SR}(\mathcal{C})}(K, L) \longrightarrow \Mor_\mathcal{C}(U, V)
$$
comme suit. Étant donné $f : K \to L$ défini par $\alpha : I \to J$
et $f_i : U_i \to V_{\alpha(i)}$, nous obtenons une transformation de foncteurs
$$
\Mor_\mathcal{C}(V, -) =
\prod\nolimits_{j \in J} \Mor_\mathcal{C}(V_j, -)
\to
\prod\nolimits_{i \in I} \Mor_\mathcal{C}(U_i, -) =
\Mor_\mathcal{C}(U, -)
$$
qui envoie $(g_j)_{j \in J}$ sur
$(g_{\alpha(i)} \circ f_i)_{i \in I}$. Le lemme de Yoneda
produit donc l'application correspondante $U \to V$. Bien sûr, $U \to V$
envoie la composante $U_i$ dans la composante $V_{\alpha(i)}$ par
le morphisme $f_i$.
\end{remark}

\begin{remark}
\label{remark-take-unions-hypercovering}
Soit $\mathcal{C}$ un site. Supposons que $\mathcal{C}$ admette
des produits fibrés et des égalisateurs, et soit $K$ un hyperrecouvrement.
Écrivons $K_n = \{U_{n, i}\}_{i \in I_n}$. Supposons que
\begin{enumerate}
\item[(a)] $U_n = \coprod_{i \in I_n} U_{n, i}$ existe ;
\item[(b)] $\coprod_{i \in I_n} h_{U_{n, i}} \to h_{U_n}$ induise
un isomorphisme après faisceautisation.
\end{enumerate}
Nous obtenons alors un autre objet simplicial $L$ de $\text{SR}(\mathcal{C})$
avec $L_n = \{U_n\}$, voir la
remarque \ref{remark-taking-disjoint-unions}.
Nous affirmons que $L$ est un hyperrecouvrement.
Pour le voir, vérifions les conditions (1), (2), (3) de la
définition \ref{definition-hypercovering-variant}.
La condition (1) résulte de (b) et de la condition (1) pour $K$.
La condition (2) s'en déduit exactement de la même manière.
La condition (3) résulte de
\begin{align*}
F((\text{cosk}_n \text{sk}_n L)_{n + 1})^\#
& =
((\text{cosk}_n \text{sk}_n F(L)^\#)_{n + 1}) \\
& =
((\text{cosk}_n \text{sk}_n F(K)^\#)_{n + 1}) \\
& =
F((\text{cosk}_n \text{sk}_n K)_{n + 1})^\#
\end{align*}
pour $n \geq 1$ ; la condition pour $K$ implique donc celle pour
$L$, exactement comme en (1) et (2).
Notons que $F$ commute aux limites connexes et que la faisceautisation est exacte,
ce qui prouve les première et dernière égalités ; l'égalité intermédiaire résulte de
$F(K)^\# = F(L)^\#$ par (b).
\end{remark}

\begin{remark}
\label{remark-take-unions-hypercovering-X}
Soit $\mathcal{C}$ un site. Soit $X \in \Ob(\mathcal{C})$.
Supposons que $\mathcal{C}$ admette des produits fibrés et soit $K$ un hyperrecouvrement de $X$.
Écrivons $K_n = \{U_{n, i}\}_{i \in I_n}$. Faisons les hypothèses suivantes :
\begin{enumerate}
\item[(a)] $U_n = \coprod_{i \in I_n} U_{n, i}$ existe ;
\item[(b)] étant donnés des morphismes
$(\alpha, f_i) : \{U_i\}_{i \in I} \to \{V_j\}_{j \in J}$ et
$(\beta, g_k) : \{W_k\}_{k \in K} \to \{V_j\}_{j \in J}$
dans $\text{SR}(\mathcal{C})$ tels que
$U = \coprod U_i$, $V = \coprod V_j$ et $W = \coprod W_j$
existent, nous avons $U \times_V W =
\coprod_{(i, j, k), \alpha(i) = j = \beta(k)} U_i \times_{V_j} W_k$,
\item[(c)] si $(\alpha, f_i) : \{U_i\}_{i \in I} \to \{V_j\}_{j \in J}$
est un recouvrement au sens de la
définition \ref{definition-covering-SR}
et si $U = \coprod U_i$ et $V = \coprod V_j$ existent,
alors le morphisme correspondant $U \to V$
de la remarque \ref{remark-taking-disjoint-unions}
est un recouvrement de $\mathcal{C}$.
\end{enumerate}
Nous obtenons alors un autre objet simplicial $L$ de $\text{SR}(\mathcal{C})$
avec $L_n = \{U_n\}$, voir la
remarque \ref{remark-taking-disjoint-unions}.
Nous affirmons que $L$ est un hyperrecouvrement de $X$.
Pour le voir, vérifions les conditions (1), (2) de la
définition \ref{definition-hypercovering}.
La condition (1) résulte de (c) et de la condition (1) pour $K$,
car la condition (1) pour $K$ dit que $K_0 = \{U_{0, i}\}_{i \in I_0}$
est un recouvrement de $\{X\}$ au sens de la
définition \ref{definition-covering-SR}.
La condition (2) résulte du fait que $\mathcal{C}/X$ admet
toutes les limites finies, donc que $\text{SR}(\mathcal{C}/X)$
les admet aussi, tandis que la condition (b) dit que la
construction consistant à « prendre les sommes disjointes » commute
à ces limites finies. Ainsi, le morphisme
$$
L_{n + 1} \longrightarrow (\text{cosk}_n \text{sk}_n L)_{n + 1}
$$
est un recouvrement, puisqu'il résulte de l'application de notre foncteur
« prendre les sommes disjointes » au morphisme
$$
K_{n + 1} \longrightarrow (\text{cosk}_n \text{sk}_n K)_{n + 1}
$$
qui est supposé être un recouvrement au sens de la
définition \ref{definition-covering-SR} par la condition (2) pour $K$.
Cela a un sens, car la propriété (b) nous assure en particulier
que, si nous partons d'un diagramme fini d'objets
semi-représentables au-dessus de $X$
pour lesquels nous pouvons prendre les sommes disjointes, alors
la limite du diagramme dans $\text{SR}(\mathcal{C}/X)$
reste un objet semi-représentable au-dessus de $X$ pour lequel
nous pouvons prendre les sommes disjointes.
\end{remark}





\begin{multicols}{2}[\section{Autres chapitres}]
\noindent
Préliminaires
\begin{enumerate}
\item \hyperref[introduction-section-phantom]{Introduction}
\item \hyperref[conventions-section-phantom]{Conventions}
\item \hyperref[sets-section-phantom]{Théorie des ensembles}
\item \hyperref[categories-section-phantom]{Catégories}
\item \hyperref[topology-section-phantom]{Topologie}
\item \hyperref[sheaves-section-phantom]{Faisceaux sur les espaces}
\item \hyperref[sites-section-phantom]{Sites et faisceaux}
\item \hyperref[stacks-section-phantom]{Champs}
\item \hyperref[fields-section-phantom]{Corps}
\item \hyperref[algebra-section-phantom]{Algèbre commutative}
\item \hyperref[brauer-section-phantom]{Groupes de Brauer}
\item \hyperref[homology-section-phantom]{Algèbre homologique}
\item \hyperref[derived-section-phantom]{Catégories dérivées}
\item \hyperref[simplicial-section-phantom]{Méthodes simpliciales}
\item \hyperref[more-algebra-section-phantom]{Compléments d'algèbre}
\item \hyperref[smoothing-section-phantom]{Lissage des morphismes d'anneaux}
\item \hyperref[modules-section-phantom]{Faisceaux de modules}
\item \hyperref[sites-modules-section-phantom]{Modules sur les sites}
\item \hyperref[injectives-section-phantom]{Objets injectifs}
\item \hyperref[cohomology-section-phantom]{Cohomologie des faisceaux}
\item \hyperref[sites-cohomology-section-phantom]{Cohomologie sur les sites}
\item \hyperref[dga-section-phantom]{Algèbre différentielle graduée}
\item \hyperref[dpa-section-phantom]{Algèbre à puissances divisées}
\item \hyperref[sdga-section-phantom]{Faisceaux différentiels gradués}
\item \hyperref[hypercovering-section-phantom]{Hyperrecouvrements}
\end{enumerate}
Schémas
\begin{enumerate}
\setcounter{enumi}{25}
\item \hyperref[schemes-section-phantom]{Schémas}
\item \hyperref[constructions-section-phantom]{Constructions de schémas}
\item \hyperref[properties-section-phantom]{Propriétés des schémas}
\item \hyperref[morphisms-section-phantom]{Morphismes de schémas}
\item \hyperref[coherent-section-phantom]{Cohomologie des schémas}
\item \hyperref[divisors-section-phantom]{Diviseurs}
\item \hyperref[limits-section-phantom]{Limites de schémas}
\item \hyperref[varieties-section-phantom]{Variétés}
\item \hyperref[topologies-section-phantom]{Topologies sur les schémas}
\item \hyperref[descent-section-phantom]{Descente}
\item \hyperref[perfect-section-phantom]{Catégories dérivées de schémas}
\item \hyperref[more-morphisms-section-phantom]{Compléments sur les morphismes}
\item \hyperref[flat-section-phantom]{Compléments sur la platitude}
\item \hyperref[groupoids-section-phantom]{Schémas en groupoïdes}
\item \hyperref[more-groupoids-section-phantom]{Compléments sur les schémas en groupoïdes}
\item \hyperref[etale-section-phantom]{Morphismes étales de schémas}
\end{enumerate}
Thèmes de théorie des schémas
\begin{enumerate}
\setcounter{enumi}{41}
\item \hyperref[chow-section-phantom]{Homologie de Chow}
\item \hyperref[intersection-section-phantom]{Théorie de l'intersection}
\item \hyperref[pic-section-phantom]{Schémas de Picard des courbes}
\item \hyperref[weil-section-phantom]{Théories cohomologiques de Weil}
\item \hyperref[adequate-section-phantom]{Modules adéquats}
\item \hyperref[dualizing-section-phantom]{Complexes dualisants}
\item \hyperref[duality-section-phantom]{Dualité pour les schémas}
\item \hyperref[discriminant-section-phantom]{Discriminants et différentes}
\item \hyperref[derham-section-phantom]{Cohomologie de de Rham}
\item \hyperref[local-cohomology-section-phantom]{Cohomologie locale}
\item \hyperref[algebraization-section-phantom]{Géométrie algébrique et formelle}
\item \hyperref[curves-section-phantom]{Courbes algébriques}
\item \hyperref[resolve-section-phantom]{Résolution des surfaces}
\item \hyperref[models-section-phantom]{Réduction semi-stable}
\item \hyperref[functors-section-phantom]{Foncteurs et morphismes}
\item \hyperref[equiv-section-phantom]{Catégories dérivées de variétés}
\item \hyperref[pione-section-phantom]{Groupes fondamentaux de schémas}
\item \hyperref[etale-cohomology-section-phantom]{Cohomologie étale}
\item \hyperref[crystalline-section-phantom]{Cohomologie cristalline}
\item \hyperref[proetale-section-phantom]{Cohomologie pro-étale}
\item \hyperref[relative-cycles-section-phantom]{Cycles relatifs}
\item \hyperref[more-etale-section-phantom]{Compléments de cohomologie étale}
\item \hyperref[trace-section-phantom]{La formule des traces}
\end{enumerate}
Espaces algébriques
\begin{enumerate}
\setcounter{enumi}{64}
\item \hyperref[spaces-section-phantom]{Espaces algébriques}
\item \hyperref[spaces-properties-section-phantom]{Propriétés des espaces algébriques}
\item \hyperref[spaces-morphisms-section-phantom]{Morphismes d'espaces algébriques}
\item \hyperref[decent-spaces-section-phantom]{Espaces algébriques décents}
\item \hyperref[spaces-cohomology-section-phantom]{Cohomologie des espaces algébriques}
\item \hyperref[spaces-limits-section-phantom]{Limites d'espaces algébriques}
\item \hyperref[spaces-divisors-section-phantom]{Diviseurs sur les espaces algébriques}
\item \hyperref[spaces-over-fields-section-phantom]{Espaces algébriques sur des corps}
\item \hyperref[spaces-topologies-section-phantom]{Topologies sur les espaces algébriques}
\item \hyperref[spaces-descent-section-phantom]{Descente et espaces algébriques}
\item \hyperref[spaces-perfect-section-phantom]{Catégories dérivées d'espaces}
\item \hyperref[spaces-more-morphisms-section-phantom]{Compléments sur les morphismes d'espaces}
\item \hyperref[spaces-flat-section-phantom]{Platitude sur les espaces algébriques}
\item \hyperref[spaces-groupoids-section-phantom]{Groupoïdes dans les espaces algébriques}
\item \hyperref[spaces-more-groupoids-section-phantom]{Compléments sur les groupoïdes dans les espaces}
\item \hyperref[bootstrap-section-phantom]{Amorçage}
\item \hyperref[spaces-pushouts-section-phantom]{Sommes amalgamées d'espaces algébriques}
\end{enumerate}
Thèmes de géométrie
\begin{enumerate}
\setcounter{enumi}{81}
\item \hyperref[spaces-chow-section-phantom]{Groupes de Chow des espaces}
\item \hyperref[groupoids-quotients-section-phantom]{Quotients de groupoïdes}
\item \hyperref[spaces-more-cohomology-section-phantom]{Compléments sur la cohomologie des espaces}
\item \hyperref[spaces-simplicial-section-phantom]{Espaces simpliciaux}
\item \hyperref[spaces-duality-section-phantom]{Dualité pour les espaces}
\item \hyperref[formal-spaces-section-phantom]{Espaces algébriques formels}
\item \hyperref[restricted-section-phantom]{Algébrisation des espaces formels}
\item \hyperref[spaces-resolve-section-phantom]{Résolution des surfaces revisitée}
\end{enumerate}
Théorie des déformations
\begin{enumerate}
\setcounter{enumi}{89}
\item \hyperref[formal-defos-section-phantom]{Théorie formelle des déformations}
\item \hyperref[defos-section-phantom]{Théorie des déformations}
\item \hyperref[cotangent-section-phantom]{Le complexe cotangent}
\item \hyperref[examples-defos-section-phantom]{Problèmes de déformation}
\end{enumerate}
Champs algébriques
\begin{enumerate}
\setcounter{enumi}{93}
\item \hyperref[algebraic-section-phantom]{Champs algébriques}
\item \hyperref[examples-stacks-section-phantom]{Exemples de champs}
\item \hyperref[stacks-sheaves-section-phantom]{Faisceaux sur les champs algébriques}
\item \hyperref[criteria-section-phantom]{Critères de représentabilité}
\item \hyperref[artin-section-phantom]{Axiomes d'Artin}
\item \hyperref[quot-section-phantom]{Espaces de Quot et de Hilbert}
\item \hyperref[stacks-properties-section-phantom]{Propriétés des champs algébriques}
\item \hyperref[stacks-morphisms-section-phantom]{Morphismes de champs algébriques}
\item \hyperref[stacks-limits-section-phantom]{Limites de champs algébriques}
\item \hyperref[stacks-cohomology-section-phantom]{Cohomologie des champs algébriques}
\item \hyperref[stacks-perfect-section-phantom]{Catégories dérivées de champs}
\item \hyperref[stacks-introduction-section-phantom]{Introduction aux champs algébriques}
\item \hyperref[stacks-more-morphisms-section-phantom]{Compléments sur les morphismes de champs}
\item \hyperref[stacks-geometry-section-phantom]{La géométrie des champs}
\end{enumerate}
Thèmes de théorie des espaces de modules
\begin{enumerate}
\setcounter{enumi}{107}
\item \hyperref[moduli-section-phantom]{Champs de modules}
\item \hyperref[moduli-curves-section-phantom]{Espaces de modules de courbes}
\end{enumerate}
Divers
\begin{enumerate}
\setcounter{enumi}{109}
\item \hyperref[examples-section-phantom]{Exemples}
\item \hyperref[exercises-section-phantom]{Exercices}
\item \hyperref[guide-section-phantom]{Guide bibliographique}
\item \hyperref[desirables-section-phantom]{Desiderata}
\item \hyperref[coding-section-phantom]{Style de programmation}
\item \hyperref[obsolete-section-phantom]{Obsolète}
\item \hyperref[fdl-section-phantom]{Licence de documentation libre GNU}
\item \hyperref[index-section-phantom]{Index généré automatiquement}
\end{enumerate}
\end{multicols}


\bibliography{my}
\bibliographystyle{amsalpha}

\end{document}
